%!TEX TS-program = pdflatex

%%%%%%%%%%%%%%%%%%%%%%%%%%%%%%%%%%%%%%%%%%%%%%%%%%%%%%%%%%%%%%%%%%%%%%%%%
%%%%%%%%%%%%%%%%%%%%%%%%%%%%%%%%%%%%%%%%%%%%%%%%%%%%%%%%%%%%%%%%%%%%%%%%%
%
% Appell-Dunkl sequences and Hurwitz-Dunkl zeta functions
% by
% A. Gil Asensi and J. L. Varona
%
% Date of last modification: 2022-11-13
%
%%%%%%%%%%%%%%%%%%%%%%%%%%%%%%%%%%%%%%%%%%%%%%%%%%%%%%%%%%%%%%%%%%%%%%%%%
%%%%%%%%%%%%%%%%%%%%%%%%%%%%%%%%%%%%%%%%%%%%%%%%%%%%%%%%%%%%%%%%%%%%%%%%%

%\documentclass[11pt,a4paper]{amsart}
\documentclass[a4paper]{amsart}

%\usepackage{amssymb}

%%% \allowdisplaybreaks is a command from amsmath that allow page breaks
%%% inside a multi-line formula (align, multline, gather, ...)
\allowdisplaybreaks[3] %%% Scale from 1 (weak allow) to 4 (complete allow).
%%% If a break can be avoided, we can use \\* instead of \\

%\usepackage[inline]{enumitem}
%\setlist[enumerate,itemize]{font=\bfseries,wide=0pt}

\usepackage{url}
\usepackage[colorlinks,linkcolor=blue,citecolor=blue]{hyperref}

%%%%%%%%%%%%%%%%%%%%%%%%%%%%
%% ONLY FOR DRAFTS:
%\usepackage[utf8]{inputenc}
%\usepackage{color}
%\newcommand{\rojo}[1]{\textcolor{red}{#1}}
%\newcommand{\azul}[1]{\textcolor{blue}{#1}}
%\newcommand{\verde}[1]{\textcolor{green}{#1}}
%%%%%%%%%%%%%%%%%%%%%%%%%%%%

\usepackage{graphicx}
\usepackage{tikz}

%%%%%%%%%%%%%%%%%%%%%%%%%%%%
%% ONLY FOR DRAFTS:
%\usepackage[notcite,notref]{showkeys}
%\usepackage{refcheck}
%\usepackage[hyperpageref]{backref} % For drafts
%\usepackage{backref}
%%%%%%%%%%%%%%%%%%%%%%%%%%%%

%\newcommand{\D}{\mathbb{D}}
\newcommand{\N}{\mathbb{N}}
\newcommand{\Z}{\mathbb{Z}}
\newcommand{\Q}{\mathbb{Q}}
\newcommand{\R}{\mathbb{R}}
\newcommand{\C}{\mathbb{C}}

\newcommand{\I}{\mathcal{I}}
\newcommand{\Ea}{E_{\alpha}}
\newcommand{\cc}{\gamma}
\newcommand{\calH}{\mathcal{H}}
\newcommand{\calO}{\mathcal{O}}
\newcommand{\calP}{\mathcal{P}}
\newcommand{\B}{\mathfrak{B}}
\newcommand{\E}{\mathfrak{E}}
\renewcommand{\epsilon}{\varepsilon}
\newcommand{\La}{\Lambda_{\alpha}}
\newcommand{\Lax}{\Lambda_{\alpha,x}}
\newcommand{\binoma}[2]{\binom{#1}{#2}_{\!\alpha}}

\renewcommand{\Re}{\operatorname{Re}}
\renewcommand{\Im}{\operatorname{Im}}

\newcommand{\Li}{\operatorname{Li}}

\renewcommand{\theenumi}{\roman{enumi}}

\newcommand{\Res}{\operatornamewithlimits{Res}}

%-----------
\theoremstyle{plain}
\newtheorem{theorem}{Theorem}[section]
\newtheorem{proposition}[theorem]{Proposition}
\newtheorem{corollary}[theorem]{Corollary}
\newtheorem{lemma}[theorem]{Lemma}

\theoremstyle{definition}
\newtheorem{definition}[theorem]{Definition}

\theoremstyle{remark}
\newtheorem*{remark}{Remark}
%-----------

\numberwithin{equation}{section}

%-----------
\title[Appell-Dunkl sequences and Hurwitz-Dunkl zeta functions]%
{Appell-Dunkl sequences\\and Hurwitz-Dunkl zeta functions}
%-----------


%-----------
\author[A. Gil Asensi]{Alejandro Gil Asensi}
\address{%
       Departamento de Matem\'aticas y Computaci\'on,
       Universidad de La Rioja,
       26006 Logro\~no, Spain}
\email{algila@unirioja.es}

\author[J.~L.~Varona]{Juan L. Varona}
\address{%
       Departamento de Matem\'aticas y Computaci\'on,
       Universidad de La Rioja,
       26006 Logro\~no, Spain}
\email{jvarona@unirioja.es}
\urladdr{http://www.unirioja.es/cu/jvarona/}
%-----------

%\date{\today}
%\date{November 13, 2022}

\renewcommand{\datename}{THIS PAPER HAS BEEN PUBLISHED IN:}
\date{\textit{J.\ Math.\ Anal.\ Appl.}\ \textbf{520} (2023), no.~1, 126870; 
DOI: \url{https://doi.org/10.1016/j.jmaa.2022.126870}}
%%% Online first:
%\date{\textit{J.\ Math.\ Anal.\ Appl.}, online first December~18, 2022, 
%doi: 10.1016/j.jmaa.2022.126870.}

\thanks{Research supported by grant PID2021-124332NB-C22 from MINECO (Spain).}
%% REALMENTE ES ESTO:
%\thanks{Research supported by grant PID2021-124332NB-C22 from 
%Ministerio de Ciencia e Innovación (Spain).}

%-----------
% 11B83 Special sequences and polynomials
% 11B68 Bernoulli and Euler numbers and polynomials
% 11M35 Hurwitz and Lerch zeta functions
% 30B40 Analytic continuation of functions of one complex variable
% 33C20 Generalized hypergeometric series, pFq
% 44A20 Integral transforms of special functions
% 42A38 Fourier and Fourier-Stieltjes transforms and other transforms of Fourier type
\subjclass[2010]{Primary 11B83; Secondary 11B68, 11M35, 30B40, 33C20, 44A20}
\keywords{Appell-Dunkl polynomials, 
Mellin transform, 
%Hurwitz zeta function, 
Hurwitz-Dunkl zeta function, 
%Lerch zeta function, 
Lerch-Dunkl zeta function, 
%Bernoulli polynomials, 
Bernoulli-Dunkl polynomials, 
%Euler polynomials, 
Euler-Dunkl polynomials}
%-----------

%%%%%%%%%%%%%%%%%%%%%%%%%%%%
%---------------------
\begin{document}
%---------------------
%%%%%%%%%%%%%%%%%%%%%%%%%%%%

%-----------
%% Class amsart Warning:
%% Abstract should precede \maketitle in AMS document classes
\begin{abstract}
Dunkl theory on the real line involves some tools such as the Dunkl derivative
\[
\Lambda_\alpha f(x) = \frac{d}{dx}f(x) + \frac{2\alpha+1}{2} \frac{f(x)-f(-x)}{x}
\]
or the Dunkl exponential $E_{\alpha}(z)$ that is defined in terms of 
the Bessel functions. 
% (also, $\mathcal{I}_\alpha(z) = {}_0F_1(\alpha+1,z^2/4)$)
Taking $\alpha=-1/2$ we get $\Lambda_{-1/2} = d/dx$ and 
$E_{-1/2}(z) = e^z$, hence, the classic derivative and exponential
are particular cases.
In recent years, some papers have generalized, in a Dunkl sense, 
number theoretic concepts such as Appell sequences, and then they are called
Appell-Dunkl sequences; in particular, the so called Bernoulli-Dunkl 
and Euler-Dunkl polynomials have been defined, among others. 
Here we generalize, also in a Dunkl sense, some Hurwitz or Lerch 
zeta functions such as $\zeta(s,x) = \sum_{n=0}^{\infty} 1/(n+x)^s$ 
and, in addition, we get properties that relate those functions,
extended to the $s$-complex plane and evaluated at negative integers $s$, 
with Bernoulli-Dunkl and Euler-Dunkl polynomials. One of the results we get 
for the ``Dunkl zeta function'' $\zeta_\alpha(s)$ is
\[
 \zeta_{\alpha}(1-s) = \Gamma(s) \cos{\left(\frac{\pi s}{2}\right)} 
    \sum_{n=1}^{\infty} \frac{1}{s_n^s}, 
    \qquad \operatorname{Re}(s)>1
\]
(where $s_n$ are the positive zeros of the 
Bessel function $J_{\alpha+1}(x)$). This equation 
provides a generalization of the reflection formula of the 
Riemann zeta function, where the function $\sum_{n=1}^{\infty} 1/s_n^s$ 
is playing a similar role as $\sum_{n=1}^{\infty} 1/n^s$.
\end{abstract}
%-----------


%%%%%%%%%%%%%%%%%%%%%%%%%%%%
%-----------
\maketitle
%-----------
%%%%%%%%%%%%%%%%%%%%%%%%%%%%

%-----------
%\tableofcontents
%-----------

%-----------
\section{Introduction}
%-----------

An Appell sequence $\{P_n(x)\}_{n=0}^{\infty}$ is a sequence of polynomials defined by a Taylor generating expansion
\begin{equation}
\label{eq:As}
  A(t) e^{xt} = \sum_{n=0}^{\infty} P_n(x) \frac{t^n}{n!},
\end{equation}
where $A(t)$ is a function analytic at $t=0$ with $A(0) \ne 0$.
Since the exponential function $e^x$ is invariant under the differential operator ${d}/{dx}$, it is easy to show that $P_n(x)$ is a polynomial of degree $n$ and $P'_n(x) = nP_{n-1}(x)$. Typical examples of Appell sequences are the Bernoulli polynomials $\{B_n(x)\}_{n=0}^{\infty}$, the Euler polynomials $\{E_n(x)\}_{n=0}^{\infty}$, or the probabilistic Hermite polynomials $\{\operatorname{He}_n(x)\}_{n=0}^{\infty}$ that are defined by taking $A(t) = \frac{te^{xt}}{e^t -1}$, $\frac{2e^{xt}}{e^t +1}$ or $e^{-t^2/2}$ respectively (a slight variation is the physicists' Hermite polynomials $\{H_n(x)\}_{n=0}^{\infty}$ defined by $e^{-t^2}e^{2xt} = \sum_{n=0}^{\infty} H_n(x) \frac{t^n}{n!}$).

The Appell sequences of polynomials have been extended in many ways. One of them consists of changing the derivative operator by operators in the context of Dunkl. In \cite{DH} and \cite{CDPV}, the derivative operator was replaced by 
\begin{equation*}
%\label{eq:La}
  \La f(x) = \frac{d}{dx}f(x) + \frac{2\alpha+1}{2}
  \left(\frac{f(x)-f(-x)}{x}\right),
\end{equation*}
where $\alpha > -1$ is a fixed parameter (see \cite{Du, Ros}); observe that the case $\alpha=-1/2$ recovers the classical case $\Lambda_{-1/2} = \frac{d}{dx}$. In that setting, an Appell-Dunkl sequence $\{P_n\}_{n=0}^{\infty}$ is
a sequence of polynomials that satisfies
\begin{equation*}
%\label{eq:LaPk}
  \La P_n(x) = \big(n+(\alpha+1/2)(1-(-1)^n)\big)P_{n-1}(x)
\end{equation*}
(instead of $\La P_n = n P_{n-1}$, the previous definition with 
a different multiplicative constant in the place of $n$ is used for 
convenience with the notation). 
The Appell-Dunkl sequences can be written as a generating expansion similar to~\eqref{eq:As}, namely
\begin{equation*}
  A(t) E_\alpha(xt) = \sum_{n=0}^{\infty} \calP_n(x) \frac{t^n}{\cc_{n,\alpha}},
\end{equation*}
for a certain function $E_\alpha$ and certain constants $\cc_{n,\alpha}$ (with $E_{-1/2} = \exp$ and $\cc_{n,-1/2}=n!$); we will see the details in Section~\ref{sec:AD}. The first Appell-Dunkl sequence of polynomials studied in the mathematical literature were the so called generalized Hermite polynomials; see~\cite{Ros}. In recent years, also the Bernoulli and the Euler polynomials (among other Appell families) have been extended to the Dunkl context; see, for instance, \cite{CDPV, CMV, DPV}. These polynomials have proved to be very useful to extend some classical properties to a more general context. For instance, the Bernoulli polynomials can be used to find the values of the series $\sum_{m=1}^{\infty} 1/m^{2k}$, and the Bernoulli-Dunkl polynomials can be used to compute the Rayleigh series $\sum_{m=1}^{\infty} 1/s_m^{2k}$, where $\{s_m\}_{m=1}^\infty$ are the positive zeros of a Bessel function (note that, essentially, the sine function is a particular case of a Bessel function, and the positive zeros of the sine are $s_m = \pi m$, $m \ge 1$, so in this case the corresponding Rayleigh series reduces to $\sum_{m=1}^{\infty} 1/m^{2k}$).

In the classical case, there is a large class of Appell sequences $\{P_n(x)\}_{n=0}^\infty$ for which there is a function $H(s,x)$, entire in $s$ for fixed $x$ with $\Re x > 0$, and satisfying $H(-n,x) = P_n(x)$ for $n = 0,1,2,\dots$ For example, in the case of Bernoulli polynomials, $H$ is essentially the Hurwitz zeta function $\zeta(s,x)$ that for $\Re(s)>1$ is defined as $\zeta(s,x) = \sum_{m=0}^{\infty} (m+x)^{-s}$, and whose analytic extension to the $s$-complex plane satisfies $-n\zeta(1-n,x) = B_n(x)$. Another well-known example is the Apostol-Bernoulli polynomials, whose corresponding function $H$ is, essentially, the Lerch transcendent function (see~\cite{Ap}). More examples can be found in~\cite{BayChi, BaySim, SuHu}. The papers \cite{NRV-2018, NRV-2019} show how this can be done, in a very general way, with the help of the Mellin transform $\int_0^\infty f(t) t^{s-1}\,dt$, and provide many additional examples. 

The aim of this paper is to show how to do it in the context of Appell-Dunkl sequences. Here, there are two important difficulties. The first one is the size of $E_\alpha(t)$ when $t \to \pm \infty$. Although $E_\alpha(t)$ is a generalization of $e^t$ to the Dunkl context, it is not true that $E_\alpha(t) \sim e^t$ when $t \to \pm \infty$, but, roughly speaking, $E_\alpha(t) \sim e^{|t|}$ (except for $\alpha=-1/2$). In the above mentioned Mellin transform, a factor $e^{-t}$ in $f(t)$ greatly contributes to the convergence of the integral; however, this does not happen with $E_\alpha(-t)$. In the second place, the classical translation $f(x) \mapsto f(x+m)$ becomes a complicate operator in the Dunkl context, and this affects the summands of type $(x+m)^{-s}$ of the classical Hurwitz zeta function, which are not so simple in the new context.

The organization of this paper is as follows. In Section~\ref{sec:AD} we give the details of the Dunkl context, and the precise definitions of the Appell-Dunkl sequences. Section~\ref{sec:Dtau} gives the details of the Dunkl translation. In Section~\ref{sec:Mellin} we give a general procedure, based on the Mellin transform, to extend an Appell-Dunkl sequence $\{\calP_n(x)\}_{n=0}^{\infty}$ to an analytic function $H(s,x)$ such that $H(-n,x) = \calP_n(x)$ (actually, it is a bit different); due to the above mentioned difficulties, this is not as general as in the classical case studied in~\cite{NRV-2018}, is not valid in the whole range of~$x$, and requires some additional hypotheses. This section also studies several particular cases of Appell-Dunkl polynomials (Bernoulli-Dunkl, Euler-Dunkl, generalized Bernoulli-Dunkl, generalized Euler-Dunkl, and generalized Hermite), giving their corresponding Hurwitz-Dunkl zeta functions. In Section~\ref{sec:prop-hurwitz} we study some additional properties of these Hurwitz-Dunkl zeta functions. In particular, we show how these functions are connected with series of type $\sum_{m=1}^{\infty} 1/j_{m,\alpha}^{s}$ (where $\{j_{m,\alpha}\}_{m=1}^\infty$ are the positive zeros of the Bessel function of order~$\alpha$), by means of some formulas that resembles Riemann's functional equation for the classical $\zeta(s)$ function: if we use $\zeta_\alpha$ to denote the function associated to de Bernoulli-Dunkl polynomials, we have
\[
 \zeta_{\alpha}(1-s) = \Gamma(s) \cos{\left(\frac{\pi s}{2}\right)} 
    \sum_{m=1}^{\infty} \frac{1}{j_{m,\alpha+1}^s}, \qquad \Re(s)>1
\]
(see the details in that section). 
In Section~\ref{sec:zetaR} we study the conection of our results with the analytic continuation to the $s$-complex plane of $Z_\alpha(s) = \sum_{m=1}^{\infty} {1}/{j_{m,\alpha}^s}$, which was studied by Hawkins \cite{Haw}.
Finally, Section~\ref{sec:proof-of-ths} includes some of the technical proofs of the results presented in Section~\ref{sec:prop-hurwitz}.

%-----------
\section{Appell-Dunkl sequences}
\label{sec:AD}
%-----------

For $\alpha > -1$, let $J_\alpha$ denote the Bessel function of order $\alpha$ and,
for complex values of the variable $z$, let
\[
  \I_\alpha(z)
  = 2^\alpha \Gamma(\alpha+1)\, \frac{J_\alpha(iz)}{(iz)^\alpha}
  = \Gamma(\alpha+1) \sum_{n=0}^\infty \frac{(z/2)^{2n}}{n!\,\Gamma(n+\alpha+1)}
  = {}_0F_1(\alpha+1,z^2/4)
\]
(the function $\I_\alpha$ is a small variation of the so-called modified Bessel function of the first kind and order~$\alpha$, usually denoted by $I_\alpha$; see~\cite{Wat} or~\cite{OlMa-NIST}).
Also, again for $z \in \C$, take % us take
\begin{equation}
\label{eq:Ea}
  \Ea(z) = \I_\alpha(z) + \frac{z}{2(\alpha+1)}\,
  \I_{\alpha+1}(z)
  = e^{z} \, {}_1F_1(\alpha+1/2,2\alpha+2,-2z).
\end{equation}

Following~\cite{Du} for $\alpha\ge-1/2$ and~\cite{Ros} for $\alpha>-1$, in the real line and with the reflection group~$\Z_2$, the Dunkl operator $\La$ is defined as
\begin{equation}
\label{eq:Duop}
  \La f(x) = \frac{d}{dx}f(x) + \frac{2\alpha+1}{2}
  \left(\frac{f(x)-f(-x)}{x}\right),
\end{equation}
where $f$ is a suitable function on $\R$. If we want to specify that the variable involved in the Dunkl operator is $x$, we will use $\Lax$.
For any $\lambda \in \C$, we have
\begin{equation}
\label{eq:LDu}
  \La \Ea(\lambda x) = \Lax \Ea(\lambda x) = \lambda \Ea(\lambda x).
\end{equation}
Let us note that, when $\alpha = -1/2$, we have $\Lambda_{-1/2} = d/dx$ and $E_{-1/2}(\lambda x) = e^{\lambda x}$.

From the definition, it is easy to check that
\[
  \Ea(z) = \sum_{n=0}^\infty \frac{z^n}{\cc_{n,\alpha}},
  \qquad
  \I_{\alpha}(x) = \sum_{n=0}^\infty \frac{x^{2n}}{\cc_{2n,\alpha}},
\]
with 
\begin{equation}
\label{eq:ccna}
\cc_{n,\alpha} =
\begin{cases} 
  2^{2k}k!\,(\alpha+1)_k, & \text{if $n=2k$},\\
  2^{2k+1}k!\,(\alpha+1)_{k+1}, & \text{if $n=2k+1$},
\end{cases}
\end{equation}
and where $(a)_n$ denotes the Pochhammer symbol
\[
  (a)_n = a(a+1)(a+2) \cdots (a+n-1) = \frac{\Gamma(a+n)}{\Gamma(a)}
\]
(with $a$ a non-negative integer); of course, $\cc_{n,-1/2} = n!$.
From~\eqref{eq:ccna}, we have
\begin{equation}
\label{eq:coccc}
  \frac{\cc_{n,\alpha}}{\cc_{n-1,\alpha}} = n+(\alpha+1/2)(1-(-1)^n) =: \theta_{n,\alpha}.
\end{equation}
We also define
\[
  \binoma{n}{j} = \frac{\cc_{n,\alpha}}{\cc_{j,\alpha}\cc_{n-j,\alpha}},
\]
which becomes the ordinary binomial coefficient in the case $\alpha=-1/2$.
To simplify the notation we sometimes write $\cc_n=\cc_{n,\alpha}$ and $\theta_n=\theta_{n,\alpha}$.
For each function $A(t)$ analytic in a neighborhood of $t=0$ and with $A(0) \ne 0$, we define an Appell-Dunkl sequence $\{P_{n}(x)\}_{n=0}^{\infty}$ by means of the generating function
\begin{equation}
\label{eq:ad}
  A(t)\Ea(xt) = \sum_{n=0}^\infty P_{n}(x) \frac{t^n}{\cc_n}
\end{equation}
(additionally to the papers \cite{AlS, Ros} cited in the introduction, Appell-Dunkl sequences have been also considered, for instance, in \cite{BG, BKI, DH}).
From this definition, it is not difficult to prove that $P_{n,\alpha}(x)$ is a polynomial of degree $n$ and, moreover, $\La P_{n}(x) = \frac{\cc_{n}}{\cc_{n-1}}\, P_{n-1}(x)$ (when $\alpha=-1/2$, this becomes the classical $P'_{n}(x) = nP_{n-1}(x)$ in the Appell sequences). 

Besides the generalized Hermite polynomials that, in the Dunkl context, were studied in~\cite{Ros},  we will use the so called Bernoulli-Dunkl polynomials, Euler-Dunkl polynomials, and their corresponding generalization with an extra parameter. 

%%%%%%%%%%%%%%%%%%%%%%%%%%%%
%% Table
%%%%%%%%%%%%%%%%%%%%%%%%%%%%
\begin{table}[!tb]
\caption{Scheme that describes the process to transform the definition 
of the classical Bernoulli and Euler polynomials into the definition 
of the Bernoulli-Dunkl and Euler-Dunkl polynomials 
(and their generalizations of order~$r$).
In the classical case, we use the ``basic'' interval $[0,1]$, 
the function $\exp$ and the factorial~$n!$; in the Dunkl case with $\alpha>-1$, 
we must use the ``basic'' interval $[-1,1]$, the function $\Ea$ and~$\cc_{n,\alpha}$.}
%\small
\newcommand{\suml}{\sum\limits}
\centering
\hspace*{-1cm}%
\begin{tabular}{ccc}
\hline
\rule[-4pt]{0pt}{14pt} & Bernoulli $\mapsto$ Bernoulli-Dunkl & Euler $\mapsto$ Euler-Dunkl 
\\
\hline
\rule{0pt}{15pt}
Classical % definitions
& $\frac{te^{xt}}{e^t-1} = \suml_{n=0}^{\infty} B_n(x) \frac{t^n}{n!}$
& $\frac{2e^{xt}}{e^t+1} = \suml_{n=0}^{\infty} E_n(x) \frac{t^n}{n!}$
\\[9pt]
$x \mapsto \frac{x+1}{2}$
& $\frac{te^{xt/2}e^{t/2}}{e^t-1} = \suml_{n=0}^{\infty} B_n(\frac{x+1}{2}) \frac{t^n}{n!}$
& $\frac{2e^{xt/2}e^{t/2}}{e^t+1} = \suml_{n=0}^{\infty} E_n(\frac{x+1}{2}) \frac{t^n}{n!}$
\\
$t \mapsto 2t$
& $\frac{2te^{xt}e^{t}}{e^{2t}-1} = \suml_{n=0}^{\infty} B_n(\frac{x+1}{2}) \frac{2^n t^n}{n!}$
& $\frac{2e^{xt}e^{t}}{e^{2t}+1} = \suml_{n=0}^{\infty} E_n(\frac{x+1}{2}) \frac{2^n t^n}{n!}$
\\[9pt]
rewrite
& $\frac{2te^{xt}}{e^{t}-e^{-t}} = \suml_{n=0}^{\infty} B_n(\frac{x+1}{2}) \frac{2^n t^n}{n!}$
& $\frac{2e^{xt}}{e^t+e^{-t}} = \suml_{n=0}^{\infty} E_n(\frac{x+1}{2}) \frac{2^n t^n}{n!}$
\\[9pt]
$\exp \mapsto \Ea$ %, $n! \mapsto \cc_{n,\alpha}$
& $\frac{2t\Ea(xt)}{\Ea(t)-\Ea(-t)} 
= \suml_{n=0}^{\infty} B_{n}^{*}(\frac{x+1}{2}) \frac{2^n t^n}{\cc_{n,\alpha}}$
& $\frac{2\Ea(xt)}{\Ea(t)+\Ea(-t)} 
= \suml_{n=0}^{\infty} E_{n}^{*}(\frac{x+1}{2}) \frac{2^n t^n}{\cc_{n,\alpha}}$
\\[9pt]
rewrite
& $\frac{2(\alpha+1)\Ea(xt)}{\I_{\alpha+1}(t)} 
= \suml_{n=0}^{\infty} B_{n}^{*}(\frac{x+1}{2}) \frac{2^n t^n}{\cc_{n,\alpha}}$
& $\frac{\Ea(xt)}{\I_{\alpha}(t)}
= \suml_{n=0}^{\infty} E_{n}^{*}(\frac{x+1}{2}) \frac{2^n t^n}{\cc_{n,\alpha}}$
\\[9pt]
Dunkl % definition
& $\frac{\Ea(xt)}{\I_{\alpha+1}(t)} 
= \suml_{n=0}^{\infty} \B_{n,\alpha}(x) \frac{t^n}{\cc_{n,\alpha}}$
& $\frac{\Ea(xt)}{\I_{\alpha}(t)}
= \suml_{n=0}^{\infty} \E_{n,\alpha}(x) \frac{t^n}{\cc_{n,\alpha}}$
\\[9pt]
Generalized % definition
& $\frac{\Ea(xt)}{(\I_{\alpha+1}(t))^r} 
= \suml_{n=0}^{\infty} \B_{n,\alpha}^{(r)}(x) \frac{t^n}{\cc_{n,\alpha}}$
& $\frac{\Ea(xt)}{(\I_{\alpha}(t))^r}
= \suml_{n=0}^{\infty} \E_{n,\alpha}^{(r)}(x) \frac{t^n}{\cc_{n,\alpha}}$
\\[9pt]
\hline
\end{tabular}%
\hspace*{-1cm}
%\caption{Scheme that describes the process to transform the definition 
%of the classical Bernoulli and Euler polynomials into the definition 
%of the Bernoulli-Dunkl and Euler-Dunkl polynomials 
%(and their generalizations of order~$r$).
%In the classical case, we use the ``basic'' interval $[0,1]$, 
%the function $\exp$ and the factorial~$n!$; in the Dunkl case with $\alpha>-1$, 
%we must use the ``basic'' interval $[-1,1]$, the function $\Ea$ and~$\cc_{n,\alpha}$.}
\label{tab:scheme}
\end{table}

%%%%%%%%%%%%%%%%%%%%%%%%%%%%
% Bernolli-Dunkl polynomials
%%%%%%%%%%%%%%%%%%%%%%%%%%%%

%-----------
\subsection{Bernoulli-Dunkl polynomials}
%-----------

Following~\cite{CDPV}, we define the Bernoulli-Dunkl polynomials $\{\B_{n,\alpha}\}_{n=0}^{\infty}$ by means of the generating function
\begin{equation}
\label{eq:defbd}
  \frac{\Ea(xt)}{\I_{\alpha+1}(t)}
  = \sum_{n=0}^\infty \frac{\B_{n,\alpha}(x)}{\cc_{n,\alpha}}\,t^n.
\end{equation}
To simplify the notation we sometimes write $\B_n=\B_{n,\alpha}$ (and $\cc_n=\cc_{n,\alpha}$).

The first few Bernoulli-Dunkl polynomials are
\begin{align*}
  \B_0(x) &= 1, \quad & \B_1(x) &= x, \\
  \B_2(x) &= x^2 - \frac{\alpha+1}{\alpha+2},\quad &
  \B_3(x) &= x^3 - x,\\
  \B_4(x) &= x^4 - 2x^2 + \frac{(\alpha+4)(\alpha+1)}{(\alpha+3)(\alpha+2)},\quad &
  \B_5(x) &= x^5 - 2\,\frac{\alpha+3}{\alpha+2}\,x^3 + \frac{\alpha+4}{\alpha+2}\,x.
\end{align*}
Some of the properties of these polynomials can be seen in~\cite{CDPV}.

Before we continue, let us explain why we use ``Bernoulli-Dunkl'' to name these polynomials.
The first reason is that
\begin{equation}
\label{eq:BtoB}
  \frac{\B_{n,-1/2}(2x-1)}{2^n} = B_n(x),
\end{equation}
where $\{B_n\}_{n=0}^{\infty}$ are the Bernoulli polynomials (for the definition and properties of the Bernoulli polynomials see, for instance,~\cite{Dil-NIST} or~\cite{DKP}). Indeed, taking into account that
\[
  E_{-1/2}(x) = e^x,\quad \I_{1/2}(x) = \frac{\sin(ix)}{ix},
\]
the relation \eqref{eq:BtoB} can be deduced substituting $x$ for $2x-1$, $t$ for $t/2$ and $\alpha$ for $-1/2$ in the definition~\eqref{eq:defbd}. 
Here, we must note that the change $x \mapsto 2x-1$ in~\eqref{eq:BtoB} is very natural, because in the reflection group $\Z_2$, which is key in the standard definition of the Dunkl operator~\eqref{eq:Duop}, the symmetry plays a important role, and thus the role of $x=0$ and $x=1$ on the classical Bernoulli polynomials must be translated to the points $-1$ and~$1$.
In fact, this is the process that is explained in Table~\ref{tab:scheme} (extracted from~\cite{CMV}) to define Bernoulli-Dunkl polynomials as an extension to the Dunkl case of the classical Bernoulli polynomials. As is shown in the table, this process can be used for other classical polynomials.

Another reason to use the name Bernoulli-Dunkl polynomials for $\B_n$ is the role that they play in certain sums involving the zeros of the Bessel functions (see~\cite{CDPV}), which is a generalization of what happens in the case $\alpha=-1/2$ with the Bernoulli polynomials. 
This will appear again later in this paper; see Corollary~\ref{cor:sum-sm}.

%%%%%%%%%%%%%%%%%%%%%%%%%%%%
% Generalized Bernoulli-Dunkl polynomials
%%%%%%%%%%%%%%%%%%%%%%%%%%%%

%-----------
\subsection{Generalized Bernoulli-Dunkl polynomials}
%-----------

In the classical case, the generalized Bernoulli polynomials of order $r$ are $\{B_{n}^{(r)}(x)\}_{n=0}^\infty$, defined by
\[
  \left(\frac{t}{e^t-1}\right)^r e^{xt} 
  = \sum_{n=0}^\infty B_{n}^{(r)}(x) \frac{t^n}{n!}.
\]
They were introduced by N{\o}rlund in 1922 (see \cite{Nor, Nor-book}).
% , and they are also known as N{\o}rlund polynomials.

When $\alpha>-1$ we can also define the generalized Bernoulli-Dunkl polynomials $\{\B_{n,\alpha}^{(r)}\}_{n=0}^{\infty}$ (or $\B_n^{(r)}$) of order $r$ by means of the generating function
\begin{equation}
\label{eq:defbdr}
  \frac{\Ea(xt)}{\I_{\alpha+1}(t)^r}
  = \sum_{n=0}^\infty \frac{\B_{n,\alpha}^{(r)}(x)}{\cc_{n,\alpha}}\,t^n.
\end{equation}
% Again, we sometimes write $\B_n^{(r)}=\B_{n,\alpha}^{(r)}$.

In this case, the generalized Bernoulli polynomials and the generalized Bernoulli-Dunkl polynomials are related by
\[
  \B_{n,-1/2}^{(r)}(2x-r) = 2^n B_n^{(r)}(x).
\]
In the recent paper \cite{ExLaMiVa} we can see how the polynomials $\B_{n,\alpha}^{(r)}$ can be used in the context of Appell-Dunkl discrete sequences, in the same way that $B_n^{(r)}$ appear in the context of Appell discrete sequences and falling factorial polynomials.

%%%%%%%%%%%%%%%%%%%%%%%%%%%%
% Euler-Dunkl polynomials
%%%%%%%%%%%%%%%%%%%%%%%%%%%%

%-----------
\subsection{Euler-Dunkl polynomials}
%-----------

We define the Euler-Dunkl polynomials $\{\E_{n,\alpha}\}_{n=0}^{\infty}$ of order $\alpha>-1$ by means of the generating function
\begin{equation*}
%\label{eq:defed}
  \frac{\Ea(xt)}{\I_{\alpha}(t)}
  = \sum_{n=0}^\infty \frac{\E_{n,\alpha}(x)}{\cc_{n,\alpha}} \,t^n.
\end{equation*}
As usual, we will sometimes denote it only by $\E_{n}$, without specifying~$\alpha$.
The first few Euler-Dunkl polynomials are
% \footnote{En \cite{CMV} hab\'ia una errata en $\E_5(x)$; aqu\'i est\'a corregida.}
\begin{align*}
  \E_0(x) &= 1, \quad & \E_1(x) &= x, \\
  \E_2(x) &= x^2-1, \quad &
  \E_3(x) &= x^3 - \frac{\alpha+2}{\alpha+1}\,x, \\
  \E_4(x) &= x^4 - 2\,\frac{\alpha+2}{\alpha+1}\,x^2 + \frac{\alpha+3}{\alpha+1},\quad &
  \E_5(x) &= x^5 - 2\,\frac{\alpha+3}{\alpha+1}\,x^3 + \frac{(\alpha+3)^2}{(\alpha+1)^2}\,x.
\end{align*}

These polynomials are related to the classical Euler polynomials $\{E_n\}_{n=0}^{\infty}$ by
\begin{equation}
\label{eq:EtoE}
  \frac{\E_{n,-1/2}(2x-1)}{2^n} = E_n(x)
\end{equation}
(for the definition and properties of the Euler polynomials see, for instance,~\cite{Dil-NIST}).
This process has been sketched in Table~\ref{tab:scheme}.

%%%%%%%%%%%%%%%%%%%%%%%%%%%%
% Generalized Euler-Dunkl polynomials
%%%%%%%%%%%%%%%%%%%%%%%%%%%%

%-----------
\subsection{Generalized Euler-Dunkl polynomials}
%-----------

When $\alpha>-1$ we can also define the generalized Euler-Dunkl polynomials $\{\E_{n,\alpha}^{(r)}\}_{n=0}^{\infty}$ (or $\E_n^{(r)}$) of order $r$ by means of the generating function
\begin{equation*}
%\label{eq:defedr}
  \frac{\Ea(xt)}{\I_{\alpha}(t)^r}
  = \sum_{n=0}^\infty \frac{\E_{n,\alpha}^{(r)}(x)}{\cc_{n,\alpha}}\,t^n.
\end{equation*}
%Again, we sometimes write $\E_n^{(r)}=\E_{n,\alpha}^{(r)}$. 

In the classical case, the generalized Euler polynomials of order $r$ are $\{E_{n}^{(r)}(x)\}_{n=0}^\infty$ defined by
\[
  \left(\frac{2}{e^t+1}\right)^r e^{xt} 
  = \sum_{n=0}^\infty E_{n}^{(r)}(x) \frac{t^n}{n!}.
\]
The generalized Euler polynomials and the generalized Euler-Dunkl polynomials are related by
\[
  \E_{n,-1/2}^{(r)}(2x-r) = 2^n E_n^{(r)}(x).
\]

%%%%%%%%%%%%%%%%%%%%%%%%%%%%
% Dunkl-traslation
%%%%%%%%%%%%%%%%%%%%%%%%%%%%

%-----------
\section{The Dunkl translation: definition and some properties}
\label{sec:Dtau}
%-----------

The Dunkl translation operator of a function $f$ is defined by
\begin{equation}
\label{eq:translation}
  \tau_y f(x) = \sum_{n=0}^\infty \La^{n}f(x) \frac{y^n}{\cc_{n,\alpha}},
  \qquad \alpha>-1,
\end{equation}
where $\La^0$ is the identity operator and $\La^{n+1} = \La(\La^n)$.  
As in the case of $\Lax = \La$, we sometimes use $\tau_{y,x}$ if we want to indicate that the translation $\tau_{y}$ is acting on a function whose variable is~$x$.
In the case $\alpha=-1/2$, the translation $\tau_y f$ is just the Taylor expansion of a function $f$ around a fixed point $x$, that is,
\[
  f(x+y) = \sum_{n=0}^\infty f^{(n)}(x) \frac{y^n}{n!}.
\]
Of course, definition~\eqref{eq:translation} is valid only for $C^{\infty}$ functions, and assuming also that the series on the right is convergent. In particular, this can be guaranteed when $f$ is a polynomial, because the operator $\La$ applied to a polynomial of degree $k$ generates a polynomial of degree $k-1$, so the series~\eqref{eq:translation} has only a finite number of nonzero summands. 
Other properties of the translation operator $\tau_y$
can be found in \cite{Ros}, \cite{Rosler-TAMS}, \cite{ThX} and~\cite{LiLiao}, 
including some integral expressions that can be applied to a wider class of functions than~\eqref{eq:translation}.

From the definition~\eqref{eq:translation}, it is clear that $\tau_y$ commutes with the Dunkl operator~$\La$.
In what follows, we are going to see some other basic properties.
It is not difficult to prove these properties, and here we state most of them without a proof; in most cases, more details can be found in~\cite{CMV}.

A nice property of the Dunkl translation, which resembles the Newton binomial $(x+y)^n = \sum_{k=0}^n \binom{n}{k} y^k x^{n-k}$, is the following:
\begin{equation}
\label{eq:Newton}
  \tau_y ((\cdot)^n)(x) = 
  \sum_{k=0}^n \binoma{n}{k} y^k x^{n-k}.
\end{equation}
More generally, and in relation to the Appell-Dunkl sequences $\{P_n(x)\}_{n=0}^{\infty}$ defined as in~\eqref{eq:ad}, the Dunkl translation satisfies 
\begin{equation*}
%\label{eq:tau-A}
  \tau_y(P_k)(x)
  = \sum_{j=0}^k \binoma{k}{j} P_j(x) y^{k-j},
\end{equation*}
which in the classical case $\alpha=-1/2$ becomes $P_k(x+y) = \sum_{j=0}^k \binom{k}{j} P_j(x) y^{k-j}$.

Another important property is the fact
\begin{equation}
\label{eq:tyxtxy}
  \tau_y f(x) = \tau_x f(y).
\end{equation}
This is a direct consequence of the above mentioned integral expressions for the Dunkl translation. 
Moreover, at least for polynomials, it can be easily checked starting from~\eqref{eq:translation} using the linearity of $\tau_y$ and its behaviour on $f(x) = x^n$, $n=0,1,2,\dots$ Indeed, using $\binoma{n}{k} = \binoma{n}{n-k}$ and~\eqref{eq:Newton} we have $\tau_y ((\cdot)^n)(x) = \tau_x ((\cdot)^n)(y)$, and this proves~\eqref{eq:tyxtxy}.

%In particular, observe that
%$\tau_y f(0) = \tau_0 f(y) = f(y)$,
%and then, by \eqref{eq:translation}, we get the power expansion
%\[
%  f(y) = \sum_{n=0}^\infty \La^{n}f(0) \frac{y^n}{\cc_{n}}.
%\]
%% (at least, for polynomials).

The inverse operator of $\tau_y$ defined as in~\eqref{eq:translation} is 
\begin{equation}
\label{eq:tau-inv}
  \tau_y^{-1} f(x) = \sum_{n=0}^\infty \frac{c_n y^n}{\cc_{n,\alpha}}\La^{n}f(x),
\end{equation}
where $c_0 = 1$ and $c_n$ for $n\ge1$ is defined by the recurrence $c_n = -\sum_{j=0}^{n-1} \binoma{n}{j} c_j$ (a proof can be found in~\cite[Lemma~4.4]{CMV}). The operator $\tau_y^{-1}$ is not, in general, a translation (in particular, it is not $\tau_{-y}$ except when $\alpha=-1/2$).

Moreover, it is not difficult to check that the operators of type $\tau_a$, $\tau_b$, $\tau_c^{-1}$ and $\tau_d^{-1}$ commute; for instance, $\tau_a \tau_b = \tau_b \tau_a$, $\tau_c^{-1} \tau_d^{-1} = \tau_d^{-1} \tau_c^{-1}$, $\tau_a \tau_c^{-1} = \tau_c^{-1} \tau_a$ and so on. Note that, in general (except when $\alpha=-1/2$), $\tau_a \tau_b$ is not a new translation, even if $a = b$. 

In relation to $\Ea$, the Dunkl translation has a nice behaviour that resembles the classical $e^{t(x+y)} = e^{tx} e^{ty}$, namely
\begin{equation}
\label{eq:tauyEa}
  \tau_y (\Ea(t \cdot))(x) = \Ea(t x)\Ea(t y).
\end{equation}
Indeed, using $\Lax (\Ea(tx)) = t \Ea(tx)$ (this is $\frac{d}{dx} e^{tx} = t e^{tx}$ 
in the classical case), the proof of \eqref{eq:tauyEa} is a simple consequence
of the definition~\eqref{eq:translation}:
\begin{align*}
  \tau_y \big(\Ea(t\cdot)\big)(x)
  &= \sum_{m=0}^{\infty} \Lambda_{\alpha,x}^m \Ea(tx) \frac{y^m}{\cc_m}
= \sum_{m=0}^{\infty} \Ea(tx) \frac{(ty)^m}{\cc_m} 
  = \Ea(tx) \Ea(ty).
\end{align*}
It is also easy to check that
\[
\tau_y^{-1} \big( \Ea(t\cdot) \big)(x)  = \Ea(tx)/ \Ea(ty).
\] 

From these relations, we can easily state the following lemmas, which we will use later in this paper:

\begin{lemma}
\label{lem:tau-n-n}
Let $\tau_y$ be the Dunkl translation operator. Then the identities
\[
\tau_y^n \big( \Ea(t\cdot) \big)(x) = \Ea(tx)\Ea(ty)^n
\]
and
\[
\tau_y^{-n} \big( \Ea(t\cdot) \big)(x)  = \Ea(tx)/ \Ea(ty)^n
\] 
holds for all $n = 0,1,2,3,\dots$
\end{lemma}

\begin{lemma}
\label{lem:tau-n-m}
Let $\tau_y$ and $\tau_z$ be Dunkl translations and let $n$ and $m$ be two non-negative integers. Then
\[
\tau_y^n \tau_z^{-m} \big( \Ea(t\cdot) \big) = \Ea(t\cdot)\Ea(ty)^n/\Ea(tz)^{m}. 
\]
\end{lemma}

There are still a couple of technical lemmas about the behaviour of the Dunkl translation which we will use later in the paper (Subsections~\ref{subsec:BD} and~\ref{subsec:ED}). We will apply these results only to functions like $\Ea$, so we can use the Dunkl translation operator \eqref{eq:translation}, which is valid only for functions in~$\mathcal{C}^{\infty}$. Then, we can assume in the lemmas and in the proofs that the functions are in~$\mathcal{C}^{\infty}$ (this could be weakened using integral expressions for the translation).

\begin{lemma}
Let $\Lambda_{\alpha,x}$ be the Dunkl operator acting over the variable $x$ and let $g(t,x)$ be a function such as the integral $\int_{0}^{\infty} g(t,x) \,dt$ converges and $\Lambda_{\alpha,x} g(t,x)$ exists. Then,
\[
  \Lambda_{\alpha,x} \int_{0}^{\infty} g(t,x) \,dt 
  = \int_{0}^{\infty} \Lambda_{\alpha,x} g(t,x) \,dt.
\]
\end{lemma}

\begin{proof}
Using the definition of $\Lambda_{\alpha,x}$, we have
\begin{align*}
  \Lambda_{\alpha,x} \int_{0}^{\infty} g(t,x) \,dt
  &= \frac{d}{dx} \int_{0}^{\infty} g(t,x) \,dt
  + \frac{2\alpha+1}{2} \frac{\int_{0}^{\infty} g(t,x) \,dt 
  - \int_{0}^{\infty} g(-t,x) \,dt}{x} \\
  &= \int_{0}^{\infty} \frac{d}{dx} g(t,x) \,dt
  + \frac{2\alpha+1}{2} \int_{0}^{\infty} \frac{g(t,x) -  g(-t,x)}{x} \,dt  \\
  & = \int_{0}^{\infty} \Lambda_{\alpha,x} g(t,x) \,dt.
\qedhere
\end{align*}
\end{proof}

\begin{lemma}
\label{lem:tau-integral}
Let $g(t,x)$ be a function in $\mathcal{C}^{\infty}$ such that the integral $\int_{0}^{\infty} g(t,x) \,dt$ converges, and let $\tau_{y,x}$ be the Dunkl translation operator. Then
\[
  \tau_{y,x} \int_{0}^{\infty} g(t,x) \,dt = \int_{0}^{\infty} \tau_{y,x} g(t,x) \,dt.
\]
\end{lemma}

\begin{proof}
By the previous lemma,
\begin{align*}
  \tau_{y,x} \int_{0}^{\infty} g(t,x) \,dt
  &= \sum_{n=0}^{\infty} \Lambda_{\alpha,x}^n \Big(\int_{0}^{\infty} g(t,x) \,dt \Big)
  \frac{y^n}{\cc_n}
  = \sum_{n=0}^{\infty} \Big(\int_{0}^{\infty} \Lambda_{\alpha,x}^n g(t,x) \,dt \Big)
  \frac{y^n}{\cc_n} \\
  &=  \int_{0}^{\infty} \sum_{n=0}^{\infty} \Lambda_{\alpha,x}^n g(t,x) \frac{y^n}{\cc_n} \,dt 
  = \int_{0}^{\infty} \tau_{y,x} g(t,x) \,dt.
\qedhere
\end{align*}
\end{proof}

%-----------
\section{The Mellin transform to get Appell-Dunkl polynomials as values of Hurwitz-Dunkl zeta functions}
\label{sec:Mellin}
%-----------

In this section, we define a special function, $H(s,x)$, which generalizes the Appell-Dunkl polynomials in such way that $H(-n,x)$ will give us the $n$-th Appell-Dunkl polynomial $P_n(x)$ multiplied by some constant. We express $H(s,x)$ in terms of the well-known Mellin transform
\[
\mathcal{M}(f)(s) =  \frac{1}{\Gamma(s)} \int_{0}^{\infty} f(t)t^{s-1} \,dt.
\]

Here we have a very general result, and later we study particular cases of generating functions that involve Bernoulli-Dunkl and Euler-Dunkl polynomials (and also their respective generalized families) and obtain particular special functions, $H(s,x)$ for each one. In Theorem~\ref{teo:mellinBD} we relate this $H(s,x)$ with a function which we call Hurwitz-Dunkl zeta function, $\zeta_{\alpha}(s,x)$ (see Definition~\ref{def:HD-zeta}), because it plays a similar role as the traditional Hurwitz zeta function $\zeta(s,x) = \sum_{n=0}^{\infty}1/(x+n)^s$ and, in addition, it generalizes $\zeta(s,x)$ when changing $\alpha=-1/2$, $x\mapsto 2x-1$ and $t\mapsto t/2$, as we explain with more detail later. On the other hand, we obtain a similar function in the Euler-Dunkl case, $\zeta_{E,\alpha}(s,x)$ (see Definition~\ref{def:HD-zeta-euler}), which generalizes the so-called Hurwitz zeta function of Euler type, $\zeta_E(s,x) = \sum_{n=0}^{\infty}(-1)^n/(x+n)^s$. 
Note that in this theorem we could assume $k=0$ (which corresponds to the usual case $A(0) \ne 0$), but we allow a more general case.

\begin{theorem}
\label{teo:MellinAD}
Let $\{\calP_{n}(x)\}_{n=0}^{\infty}$ be an Appell-Dunkl sequence with generating function $G(x,t) = A(t) \Ea(xt)$ and suppose that $A(t)$ has a zero of order $k$ at $t=0$. 
We also assume that, for all $x \in (a,b)$, the integral
\begin{equation}
\label{eq:MellinAD}
    H(s,x) = \frac{1}{\Gamma(s)}\int_{0}^{\infty} G(x,-t)\,t^{s-1}\,dt 
         = \frac{1}{\Gamma(s)} \int_{0}^{\infty} A(-t) \Ea(-xt) t^{s-1}\,dt
\end{equation}
converges in the right plane $\Re(s) > -k$ to a holomophic function.
Then, $H(s,x)$ may be analytically continued to an entire function of $s$ satisfying
\begin{equation*}
%\label{E:sinCeros}
  H(-n,x) = \frac{n!}{\gamma_{n,\alpha}}\calP_{n}(x),\quad n=0,1,2,\dots
\end{equation*}
\end{theorem}

\begin{proof}
Suppose $H(s,x)$ converges in the right plane $\Re(s) > -k$ for all $x \in (a,b)$, as was stated in the hypothesis of the theorem.
Given $N \in \N \cup \{0\}$ with $N \geq k$, the Mellin integral can be analytically continued to the half plane $\Re(s) > -N-1$ as follows. Fix $r$ with $0 < r < R$ and $x$ with $a < x < b$ and separate the complete integral into three parts:
\begin{align*}
  H(s,x) &= \frac{1}{\Gamma(s)} \int_r^\infty A(-t) \Ea(-xt)t^{s-1}\,dt \\
  & \qquad + \frac{1}{\Gamma(s)} \int_0^{r} \left( A(-t) \Ea(-xt)
  - \sum_{n=0}^N \calP_{n}(x)\frac{(-t)^n}{\gamma_{n,\alpha}} \right)t^{s-1}\,dt \\
  & \qquad + \frac{1}{\Gamma(s)} \int_0^r \sum_{n=0}^N \calP_{n}(x) 
  \frac{(-t)^n}{\gamma_{n,\alpha}}t^{s-1} \,dt.
\end{align*}
In the first part, the integrand is $\Ea(-xt)A(-t)t^{s-1}$. Since $a<x<b$, it converges when $t \to \infty$, hence the integral is an entire function of $s$, dominated on arbitrary closed vertical strips of finite width. We may conclude that the integral is an entire function of~$s$.

In the second part, the integrand is the product of $t^{s-1}$ with the tail of the generating series, $\sum_{n=N+1}^{\infty} \calP_{n}(x) (-t)^{n}/\gamma_{n,\alpha}$, which, since $|t| \leq r < R$, is $\calO(t^{N+1})$ at $t=0$. Thus, for $\Re(s) > -N-1$, the complete integrand is $\calO(t^{N+\Re(s)})$ at $t=0$ (with the order constant depending only on $x$) and hence is integrable on $[0,r]$ and dominated on closed vertical sub-strips of finite width of this section of the $s$-plane. Therefore the second integral is a holomorphic function of $s$ for $\Re(s) > -N-1$.

In the third part, we have
\begin{align*}
\frac{1}{\Gamma(s)} \int_0^r \sum_{n=0}^N 
\calP_{n}(x) \frac{(-t)^n}{\gamma_{n,\alpha}}t^{s-1} \,dt 
& = \frac{1}{\Gamma(s)} \sum_{n=0}^N \calP_{n}(x) \frac{(-1)^n}{\gamma_{n,\alpha}}
\int_0^r t^{s+n-1} \,dt \\ 
& = \frac{1}{\Gamma(s)}\sum_{n=0}^N \calP_{n}(x) 
\frac{(-1)^n}{\gamma_{n,\alpha}} \frac{r^{s+n}}{s+n},
\end{align*}
which is an entire function of $s$ because of the simple pole of $\Gamma(s)$ at $s = -n$ cancels the simple zero of $s+n$ for $n=0,1,2,\dots$, leaving the non-zero residue $(-1)^n/n!$. 

Finally, if $s = -n$ with $0 \leq n \leq N$, the $1/\Gamma(s)$ factors in front of the first two terms vanish, as well as every term in the sum except the one corresponding to $n$, where the remaining value is $\calP_{n}(x)n!/\gamma_{n,\alpha}$ because of the residue of $\Gamma(s)$ at $-n = 0,1,2,\dots$ Thus $H(-n,x) = \calP_{n}(x)n!/\gamma_{n,\alpha}$ for these $n$ and, as $N \geq k$ was arbitrary, this completes the proof.
\end{proof}

The previous theorem is very general but needs the convergence of~\eqref{eq:MellinAD}. In the classical case $\alpha=-1/2$ stated in~\cite{NRV-2018}, we have $E_{-1/2}(-t) = e^{-t}$, which tends very quickly to $0$ when $t\to\infty$. This allows us to prove the convergence of 
\eqref{eq:MellinAD} with very weak hypothesis for $A(t)$. 
For instance, when $G(x,t)$ is the generating function for the Bernoulli polynomials,
\eqref{eq:MellinAD} becomes
\[
   H(s,x) = \frac{1}{\Gamma(s)} \int_{0}^{\infty} G(x,-t)\,t^{s-1}\,dt 
         = \frac{1}{\Gamma(s)} \int_{0}^{\infty} \frac{t}{1-e^{-t}} e^{-xt} t^{s-1}\,dt
\]
and it is clear that this integral is convergent for every $x>0$ when $s$ is in right half-plane $\Re(s) > 0$.

But this is no longer true when $\alpha \ne -1/2$.
Actually, let us recall that $\Ea(z) = e^{z} \, {}_1F_1(\alpha+1/2,2\alpha+2,-2z)$.
For proving the convergence of the integral we need to estimate the size of the integrand in~\eqref{eq:MellinAD}; in particular, the size of the factor $\Ea(-xt)$.

With this aim, let us use the asymptotic expansions of the Kummer confluent hypergeometric function ${}_1F_1(\cdot , \cdot , z)$ for $|z| \to \infty$ in the sectors
\begin{gather*}
S_+ = \left\lbrace z \in \C: -\pi/2 < \arg(z) < 3\pi/2 \right\rbrace , \\
S_- = \left\lbrace z \in \C: -3\pi/2 < \arg(z) < \pi/2 \right\rbrace 
\end{gather*}
(see, for instance \cite[p.~128]{Rosler-LN}).
In our case, these asymptotic expansions are, respectively, of the form
\begin{equation}
\label{eq:asymptotic-1f1}
\begin{aligned}
& {}_1F_1\left(\frac{2\alpha+1}{2},2\alpha+2,z\right) 
= \frac{\Gamma(2\alpha+2)}{\Gamma(\frac{2\alpha+1}{2})}e^z z^{-\alpha-3/2} 
\left( 1 + \calO\Big(\frac{1}{|z|}\Big)\right) \\
&\qquad\qquad\qquad\qquad 
+ \frac{\Gamma(2\alpha+2)}{\Gamma(\frac{2\alpha+3}{2})} e^{\pm{(2\alpha+1)i\pi}/{2}} 
z^{-\alpha-1/2} \left( 1 + \calO\Big(\frac{1}{|z|}\Big)\right) .
\end{aligned}
\end{equation}
Notice that, in the case $\alpha=-1/2$, the coefficient of the first summand is $\Gamma(1)/\Gamma(0) = 0$, so the first summand vanishes. Otherwise (for simplicity, let us assume here that the variable $z$ is real), the ``exponential parts'' for $\Ea(z) = e^{z} \, {}_1F_1(\alpha+1/2,2\alpha+2,-2z)$ in \eqref{eq:asymptotic-1f1} appears as $e^{-z}$ in the first summand, and as $e^{z}$ in the second summand. Then, the asymptotic size $e^{-t}$ (for $t\to\infty$) of the classical case $\alpha=-1/2$ becomes something similar to $\Ea(-t) \sim e^{|t|}$ for $\alpha\ne-1/2$. In this way, instead of ``a help'' to prove the convergence of~\eqref{eq:MellinAD}, the factor $\Ea(-xt)$ is a handicap, and a further analysis will be necessary to state the convergence of~\eqref{eq:MellinAD}. 

On the other hand, we would like to rewrite the function $H(s,x)$ that appears in Theorem~\ref{teo:MellinAD} as a series, just as it occurs in the classical zeta function.

For the generating function of the Bernoulli polynomials, the Mellin transform of the $G(x,-t)$ is, for $x>0$ and $\Re(s)>0$,
\begin{equation}
\label{eq:trickGeo}
\begin{aligned}
  H(s,x) &= \frac{1}{\Gamma(s)} \int_{0}^{\infty} G(x,-t)\,e^{-xt}\,t^{s-1}\,dt 
  = \frac{1}{\Gamma(s)} \int_{0}^{\infty} \frac{t}{1-e^{-t}}\,e^{-xt}\,t^{s-1}\,dt 
  \\
  &= \frac{1}{\Gamma(s)} \int_{0}^{\infty} \sum_{n=0}^\infty e^{-nt} \,e^{-xt}\,t^{s}\,dt
  = \frac{1}{\Gamma(s)} \sum_{n=0}^\infty \int_{0}^{\infty} e^{-(n + x)t}\,t^{s}\,dt
  \\
  &= \frac{\Gamma(s+1)}{\Gamma(s)} \sum_{n=0}^\infty \frac{1}{(n + x)^{s+1}}
  = s \sum_{n=0}^\infty \frac{1}{(n + x)^{s+1}}
\end{aligned}
\end{equation}
(a similar method can be followed for the Euler polynomials, as well as other Appell sequences; see, for instance, \cite{NRV-2018}).
Then, $H(s,x)$ can be given, for $x>0$ and $\Re(s)>0$, in terms of the Hurwitz zeta function
\begin{equation}
\label{eq:zetaserie}
\zeta(s,x) = \sum_{n=0}^{\infty} \frac{1}{(x+n)^{s}},
\end{equation}
so $H(s,x)$ is the analytic continuation (in the variable $s$) of $s\zeta(s+1,x)$ for $\zeta$ defined in~\eqref{eq:zetaserie} by means of a series that converges in a certain domain. Or, with the same meaning, we can say that the analytic continuation of the function $\zeta(s,x)$ defined in~\eqref{eq:zetaserie} is $H(s-1,x)/(s-1)$. 

Let us finally note that, in the previous example related to the Bernoulli polynomials, 
$A(-t)$ was written, essentially, as a geometric series 
$\sum_{n=0}^{\infty} {(e^{-t})}^n$, and then $H(s,x)$ was computed as 
a series where there was a way to compute each summand.
The analogous behavior for the Dunkl case is much more cumbersome. 
Not only is it not possible to express the integrals by means of well-known standard functions, but also the summands $1/(x+n)$ of the series become Dunkl translations instead of ordinary translations. 

%-----------
\subsection{The Bernoulli-Dunkl case}
\label{subsec:BD}
%-----------

To adapt Theorem~\ref{teo:MellinAD} to the case of the Bernoulli-Dunkl polynomials defined in~\eqref{eq:defbd}, let us first note that $\I_\alpha(t)$ is an even function, so the denominator in the left hand side of \eqref{eq:defbd} can be written as 
\begin{equation}
\label{eq:IEE}
\I_{\alpha+1}(t) = \frac{\alpha+1}{t}(\Ea(t)-\Ea(-t)).
\end{equation}

Then, concerning \eqref{eq:MellinAD} for the Bernoulli-Dunkl case, we have the following:

\begin{lemma}
\label{lem:Mellin-BD}
For $\alpha>-1$ and $x \in (-1,1)$, the integral
\[
    H(s,x) = \frac{1}{\Gamma(s)} 
    \int_{0}^{\infty} \frac{\Ea(-xt)}{\I_{\alpha+1}(t)} t^{s-1}\,dt
    = \frac{1}{(\alpha+1)\Gamma(s)}
    \int_{0}^{\infty} \frac{\Ea(-xt)}{\Ea(t)-\Ea(-t)} t^{s}\,dt
\]
converges in the right plane $\Re(s) > 0$ to a holomorphic function.
\end{lemma}

\begin{proof}
The convergence of the integral for $t$ near $0$ is clear, 
so let us analyse what happens when $t \to \infty$.
By using \eqref{eq:asymptotic-1f1}, we have, for $|z|\to\infty$,
%% Es m\'as c\'omodo usar $\Ea(-z)$ para no tener que arrastrar 
%% los $-$ que saldr\'ian si us\'aramos $\Ea(z)$
\begin{equation*}
%\label{eq:asymptotic-Ea}
\begin{aligned}
\Ea(-z) &= e^{-z} \,{}_1F_1\left(\frac{2\alpha+1}{2},2\alpha+2,2z\right) \\
&= \frac{\Gamma(2\alpha+2)}{\Gamma(\frac{2\alpha+1}{2})} e^z (2z)^{-\alpha-3/2} 
\left( 1 + \calO\Big(\frac{1}{|z|}\Big)\right) \\
&\qquad\quad 
+ \frac{\Gamma(2\alpha+2)}{\Gamma(\frac{2\alpha+3}{2})} e^{\pm{(2\alpha+1)i\pi}/{2}} 
e^{-z} (2z)^{-\alpha-1/2} \left( 1 + \calO\Big(\frac{1}{|z|}\Big)\right) .
\end{aligned}
\end{equation*}
For simplicity, let us write it as
\[
\Ea(-z) = C_1 e^z z^{-\alpha-3/2} \left( 1 + \calO\Big(\frac{1}{|z|}\Big)\right)
+ C_2^{\pm} e^{-z} z^{-\alpha-1/2} \left( 1 + \calO\Big(\frac{1}{|z|}\Big)\right).
\]
When $\alpha=-1/2$ then $C_1 = 0$; but this case is well-known and we do not need to analyse it. Then, let us assume that $\alpha \ne -1/2$.

Now, let us suppose that $x > 0$. Then we have, for $t \to \infty$
(without lost of generality we can assume $|xt|>1$), 
\begin{equation}
\label{asint-c3}
\Big|\frac{\Ea(-xt)}{\Ea(t)-\Ea(-t)} t^{s} \Big| = 
C_3 \frac{e^{xt} (xt)^{-\alpha-3/2} t^{\Re(s)} ( 1 + \calO(|xt|^{-1})) }
{e^t t^{-\alpha-1/2} ( 1 + \calO(|t|^{-1}))};
\end{equation}
this guarantees the convergence of the integral for $0 \le x < 1$. 
For $x < 0$ we have
\begin{equation}
\label{asint-c4}
\Big|\frac{\Ea(-xt)}{\Ea(t)-\Ea(-t)} t^{s} \Big| = 
C_4 \frac{e^{|xt|} |xt|^{-\alpha-1/2} t^{\Re(s)} ( 1 + \calO(|xt|^{-1})) }
{e^t t^{-\alpha-1/2} ( 1 + \calO(|t|^{-1}))},
\end{equation}
and this guarantees the convergence of the integral for $-1 < x \le 0$. 

By standard arguments on differentiation of parametric integrals, together with the above estimates, the function $H(s,x)$ is holomorphic on~$s$.
\end{proof}

The above lemma proves the hypothesis of Theorem~\ref{teo:MellinAD} for $x\in(-1,1)$. Then, we have the following:

\begin{theorem}
\label{teo:mellinBD}
Let $\Ea(xt)/\I_{\alpha+1}(t)$ be the generating function of Bernoulli-Dunkl polynomials. 
Then for $x \in (-1,1)$, the integral
\begin{equation}
\label{eq:mellinBD}
H(s,x) = \frac{1}{\Gamma(s)} 
\int_{0}^{\infty} \frac{\Ea(-xt)}{\I_{\alpha+1}(-t)} t^{s-1}\,dt
= \frac{1}{(\alpha+1)\Gamma(s)}
\int_{0}^{\infty} \frac{\Ea(-xt)}{\Ea(t)-\Ea(-t)} t^{s}\,dt
\end{equation}
converges in the right plane $\Re(s) > 0$ to a holomorphic function, which may be analytically continued to an entire function of $s$ satisfying
\begin{equation*}
%\label{eq:H-nx}
  H(-n,x) = \frac{n!}{\gamma_{n,\alpha}}\B_{n,\alpha}(x),\quad n=0,1,2,\dots
\end{equation*}
\end{theorem}

The next step is to try to write $H(s,x)$, for $\Re(s)>0$, as a kind of Hurwitz function similar to \eqref{eq:zetaserie}, as in the classical Bernoulli case.
% Let us firstly observe what happens in the classical Bernoulli case, 
% and then we will try to adapt it to the Dunkl context.

In order to compute $H(s,x)$ we may write $A(t) = 1/\I_{\alpha+1}(t)$ as a geometric series. To do that, we use the fact that $\I_{\alpha+1}(t)$ is an even function, and we use the definition of $\Ea(t)$. By \eqref{eq:IEE} we have that 
%\begin{equation*}
%\Ea(t) = \I_{\alpha}(t) + \frac{t}{2(\alpha+1)}\I_{\alpha+1}(t), 
%\quad
%\Ea(-t) = \I_{\alpha}(t) - \frac{t}{2(\alpha+1)}\I_{\alpha+1}(t),
%\end{equation*}
%so
%\[
%\I_{\alpha+1}(t) = \frac{\alpha+1}{t}\left( \Ea(t) - \Ea(-t)\right),
%\]
%and hence,
\[
A(t) = \frac{1}{\I_{\alpha+1}(t)} 
= \frac{t}{\alpha+1} \frac{1}{\Ea(t)} \frac{1}{1 - \frac{\Ea(-t)}{\Ea(t)}}
= \frac{t}{\alpha+1} \frac{1}{\Ea(t)} \sum_{n=0}^{\infty} \left(\frac{\Ea(-t)}{\Ea(t)}\right)^n.
\]
This is valid for all $t \geq 0$ and it is enough for our purposes since we just need convergence for $t \in [0, \infty)$.

Finally, we have
\begin{equation}
\label{eq:sumFBD}
\begin{aligned}
H(s,x) & = \frac{1}{\Gamma(s)} \int_{0}^{\infty} 
\frac{\Ea(-xt)}{\I_{\alpha+1}(t)} t^{s-1} \,dt \\
& = \frac{1}{\Gamma(s)} \frac{1}{\alpha+1}
\sum_{n=0}^{\infty} \int_{0}^{\infty} \frac{\Ea(-xt)}{\Ea(t)}
\left(\frac{\Ea(-t)}{\Ea(t)}\right)^n t^{s} \,dt.
\end{aligned}
\end{equation}
Notice that, in a similar way to the proof of Lemma~\ref{lem:Mellin-BD}, we can easily check that all the integrals in \eqref{eq:sumFBD} are convergent for $x \in (-1,1)$, and the interchange of the sum and the integral is justified.
However, \eqref{eq:sumFBD} is more complicated than \eqref{eq:trickGeo}; the integrals cannot be written in a closed form and we don't obtain something as simple as~\eqref{eq:zetaserie}.

In any case, we can define a kind of Hurwitz function related to the Bernoulli-Dunkl case in the following way (observe that, with the notation of \eqref{eq:sumFBD} and~\eqref{eq:sumFBD}, now we are changing $s$ to $s-1$):

\begin{definition}
\label{def:HD-zeta}
For $x \in (-1, 1)$ and $\Re(s) > 1$, we define the Hurwitz-Dunkl zeta function as
\begin{equation}
\label{eq:HD}
\zeta_{\alpha}(s,x) = \frac{1}{\Gamma(s)} \sum_{n=0}^{\infty}
\int_{0}^{\infty} \frac{\Ea(-xt)}{\Ea(t)} \left(\frac{\Ea(-t)}{\Ea(t)}\right)^n t^{s-1} \,dt.
\end{equation}
Also, we call
\begin{equation}
\label{eq:sumandoHD}
d_{\alpha}(s,x) = \frac{1}{\Gamma(s)} \int_{0}^{\infty} \frac{\Ea(-xt)}{\Ea(t)} t^{s-1} \,dt
\end{equation}
the basic Hurwitz-Dunkl term.
\end{definition}

Then we have, for $x \in (-1, 1)$ and $\Re(s)>0$,
\begin{equation}
\label{eq:HtozetaBD}
H(s,x) = \frac{s}{\alpha+1} \,\zeta_{\alpha}(s+1,x),
\end{equation}
so we can say that the function $H(s,x)$ of Theorem~\ref{teo:MellinAD} (which exists for $s\in\C$) is the analytic extension to the $s$-complex plane 
of the function $\frac{s}{\alpha+1} \,\zeta_{\alpha}(s+1,x)$; 
equivalently, we can define the analytic extension of
\begin{equation}
\label{eq:HD-zeta1}
  \zeta_{\alpha}(s,x) =  \frac{1}{\Gamma(s)}
\int_{0}^{\infty} \frac{\Ea(-xt)}{\Ea(t)-\Ea(-t)} t^{s-1}\,dt, 
\quad \Re(s) > 1
\end{equation}
(which corresponds to~\eqref{eq:HD}) as
\begin{equation}
\label{eq:HD-zeta}
  \zeta_{\alpha}(s,x) = \frac{\alpha+1}{s-1} H(s-1,x),
  \quad s \in \C,
\end{equation}
valid for $x \in (-1,1)$. 

Finally, let us see how it is possible to give an expression for $\zeta_{\alpha}(s,x)$ (valid in the half plane $\Re(s)>1$) which, in some sense, is very similar to the series~\eqref{eq:zetaserie} for the classical Hurwitz zeta function $\zeta(s,x)$, where we have a series of summands translated by means of $x \mapsto x+n$. In the Dunkl case, we are going to find an expression for $\zeta_{\alpha}(s,x)$ that, in the place of classical translations, use the Dunkl transform defined in~\eqref{eq:translation}. 

The following theorem provides an expression of the Hurwitz-Dunkl zeta function by using Dunkl translations. For simplicity, we have defined, here and in what follows, a ``symmetric translation'' $\sigma_1$ as 
\[
\sigma_1 = \tau_{1}\tau_{-1}^{-1}
\]
(or $\sigma_{1,x} = \tau_{1,x}\tau_{-1,x}^{-1}$ to clarify that it is applied to the variable~$x$). Notice that the composition of translation operators is commutative, and also the composition with inverse translations (see its expression in~\eqref{eq:tau-inv}), so we can use $\sigma_1^n = \tau_{1}^n\tau_{-1}^{-n}$ without paying attention to the order of the operators.

\begin{theorem}
\label{teo:HD-traslations}
For $\Re(s) > 1$, the Hurwitz-Dunkl zeta function can be written as
\begin{equation}
\label{eq:HD-traslations}
\zeta_{\alpha}(s,x) = \frac{1}{\Gamma(s)} \sum_{n=0}^{\infty}
\sigma_1^n \int_{0}^{\infty} \frac{\Ea(-xt)}{\Ea(t)} t^{s-1} \,dt
= \sum_{n=0}^{\infty} \sigma_1^n  d_{\alpha}(s,x).
\end{equation}
\end{theorem}

\begin{proof}[Proof of Theorem~\ref{teo:HD-traslations}]
From Lemma~\ref{lem:tau-n-n}, we have that
\begin{align*}
\tau_{1}^n \big( \Ea(-t\cdot) \big)(x) & = \Ea(-xt)\Ea(-t)^n, \\
\tau_{-1}^{-(n+1)} \big( \Ea(-t\cdot) \big)(x) & = \Ea(-xt)/ \Ea(t)^{n+1}
\end{align*}
hold. This can be easily proved as it was stated in Lemma~\ref{lem:tau-n-n} by changing $\Ea(t\cdot)$ for $\Ea(-t\cdot)$. So, from Lemma~\ref{lem:tau-n-m} we can conclude that
\begin{align*}
\tau_{1}^n \tau_{-1}^{-(n+1)} \big( \Ea(t\cdot) \big)(-x) = 
\tau_{1}^n \left( \frac{\Ea(-t\cdot)}{\Ea(t)^{n+1}} \right) (x) 
& = \Ea(-xt)\frac{\Ea(-t)^n}{\Ea(t)^{n+1}}   \\
& = \frac{\Ea(-xt)}{\Ea(t)}\left(\frac{\Ea(-t)}{\Ea(t)}\right)^n,
\end{align*}
and Lemma~\ref{lem:tau-integral} gives us that
\begin{align*}
%\label{eq:integrand}
H(s,x) & = \frac{1}{\Gamma(s)} \frac{1}{\alpha+1}
\sum_{n=0}^{\infty} \int_{0}^{\infty} \tau_{1}^n \tau_{-1}^{-(n+1)}
\big( \Ea(-t\cdot) \big)(x) t^{s} \,dt \\
&  = \frac{1}{\Gamma(s)} \frac{1}{\alpha+1}
\sum_{n=0}^{\infty} \sigma_1^n  
\int_{0}^{\infty} \tau_{-1}^{-1} \big( \Ea(-t\cdot) \big)(x) t^{s} \,dt  \\
&  = \frac{1}{\Gamma(s)} \frac{1}{\alpha+1} 
\sum_{n=0}^{\infty} \sigma_1^n 
\int_{0}^{\infty} \frac{\Ea(-xt)}{\Ea(t)} t^{s} \,dt,
\end{align*}
and the proof is concluded.
\end{proof}

Let us see that the role of $\zeta_{\alpha}(s,x)$ with the Mellin transform of Appell-Dunkl sequences is the same as the role of $\zeta(s,x)$ with the Mellin transform of Appell sequences. In fact, it generalizes the traditional Hurwitz zeta function. To see that, we observe that to transform Bernoulli polynomials into Bernoulli-Dunkl polynomials, we had to change $x \mapsto (x+1)/2$ and $t \mapsto 2t$. For that, we need to undo the change to recover the classical Hurwitz zeta function, that means, to take $\alpha=-1/2$, $x\mapsto 2x-1$ and $t\mapsto t/2$ (although many times we will not change~$t$). 

Now, let $\alpha=-1/2$. Then,
\[
d_{-1/2}(s,x) = \frac{1}{\Gamma(s)} \int_{0}^{\infty} e^{-t(x+1)}t^{s-1} \,dt 
= \frac{1}{(x+1)^{s}}
\]
if $x \in (-1,1)$. In this case, $\tau_{1}^n(f)(x) = f(x+n)$ and $\tau_{-1}^{-n}(f)(x) = f(x+n)$ and hence, 
\[
\zeta_{-1/2}(s,x) = \sum_{n=0}^{\infty} \frac{1}{(x+1+2n)^{s}}.
\]
And finally, as we are considering Bernoulli-Dunkl polynomials, we need to change $x \mapsto 2x - 1$. Hence,
\[
\zeta_{-1/2}(s,2x-1) = \sum_{n=0}^{\infty} \frac{1}{(2x+1-1+2n)^{s}} =
\frac{1}{2^s}
\sum_{n=0}^{\infty} \frac{1}{(x+n)^{s}} =  \frac{1}{2^s}\zeta(s,x).
\]
Furthermore, \eqref{eq:HD-zeta} is an integral representation of $\zeta_\alpha(s,x)$ which, as expected, generalizes under these changes the classical integral representation of~$\zeta(s,x)$:
\begin{equation*}
\zeta(s,x) = \frac{1}{\Gamma(s)}\int_{0}^{\infty} \frac{e^{-xt}}{1-e^{-t}}t^{s-1}\,dt.
\end{equation*} 
Next we summarize some of the properties of $\zeta_\alpha(s,x)$ that generalize the ones for $\zeta(s,x)$; we have already proved most of them in the preceding sections, or they are direct consequences:

\begin{proposition}[Properties of $\zeta_\alpha(s,x)$]
\label{prop:basic-z}
For $\alpha>-1$ and $x \in (-1,1)$, the function $\zeta_\alpha(s,x)$ satisfies the following:
\begin{enumerate}
\item Recurrence identities: let $\sigma_1 = \tau_1 \tau_{-1}^{-1}$; then for $\Re(s)>0$,
\begin{equation}
\label{eq:rec-HD}
    \zeta_{\alpha}(s,x) - \sigma_1(\zeta_{\alpha}(s,\cdot))(x) = d_{\alpha}(s,x),
\end{equation}
\begin{equation}
\label{eq:rec-m-HD}
    \zeta_{\alpha}(s,x) - \sigma_1^m(\zeta_{\alpha}(s,\cdot))(x) 
    = \sum_{n=0}^{m-1} \sigma_1^{n}d_{\alpha}(s,x).
\end{equation}
\item The Dunkl derivative of $\zeta_{\alpha}$:
\begin{equation}
\label{eq:prop-der-dunkl}
\Lambda_{\alpha,x} (\zeta_{\alpha}(s,x)) = -s \zeta_{\alpha}(s+1,x).
\end{equation}
\item Relation of $\zeta_{\alpha}$ with Bernoulli-Dunkl polynomials:
for $n=0,1,2,\dots$, we have
\begin{equation}
\label{eq:HLD-ber-dunkl}
    \zeta_{\alpha}(-n,x) = -\B_{n+1,\alpha}(x) \frac{(\alpha+1)n!}{\gamma_{n+1,\alpha}}.
\end{equation}
\end{enumerate}
\end{proposition}


Now we show that when $\alpha=-1/2$ and $x\mapsto 2x - 1$ we get the corresponding properties of the classical $\zeta(s,x)$. First, for the recurrence identities, we have $\sigma_1^n (f(x)) = f(x+2n)$ so $\sigma_1^n (\zeta_{\alpha}(s,\cdot))(x) = \zeta_{\alpha}(s,x+2n)$ and hence, \eqref{eq:rec-HD} transforms into 
\begin{equation*}
%\label{eq:trad-rec-HD}
    \zeta(s,x) = \zeta(s,x+1) + x^{-s}
\end{equation*}
and \eqref{eq:rec-m-HD} transforms into
\begin{equation*}
%\label{eq:trad-rec-m-HD}
    \zeta(s,x) = \zeta(s,x+m) + \sum_{n=0}^{m-1}(x+n)^{-s}
\end{equation*}
(see, for instance \cite[25.11.3 and 25.11.4]{OlMa-NIST}). Basically, the Dunkl translation $\sigma_1$ is playing the role of $x+1$ in the Hurwitz function.

In the case $\alpha=-1/2$, we have $\Lambda_{-1/2,x} = d/dx$, so \eqref{eq:prop-der-dunkl} transforms into (see, for instance, \cite[25.11.17]{OlMa-NIST})
\[
\frac{d}{dx}\zeta(s,x) = -s \zeta(s+1,x).
\]


We also get the classical relation with Bernoulli polynomials since
\begin{equation*}
\zeta_{-1/2}(-n,2x-1) = -\B_{n+1,-1/2}(2x-1)\frac{n!}{(n+1)!}\frac{1}{2} 
= -B_{n+1}(x)2^{n+1}\frac{1}{n+1}\frac{1}{2}.
\end{equation*}
Since $\zeta_{-1/2}(-n,2x-1)=2^{n}\zeta(-n,x)$, we get (see \cite[25.11.14]{OlMa-NIST})
\begin{equation}
\label{eq:bern-hurw}
    \zeta(-n,x) = -\frac{B_{n+1}(x)}{n+1}.
\end{equation}

%%%%%%%%%%%%%%%%%%%%%%%%%%%%
% The Euler-Dunkl case
%%%%%%%%%%%%%%%%%%%%%%%%%%%%

%-----------
\subsection{The Euler-Dunkl case}
\label{subsec:ED}
%-----------

This is similar to the Bernoulli-Dunkl case, but with $A(t) = 1/\I_{\alpha}(t)$.

\begin{theorem}
\label{teo:mellinED}
Let $\Ea(xt)/\I_{\alpha}(t)$ be the generating function of Euler-Dunkl polynomials.
Then for $x \in (-1,1)$, the integral
\begin{equation*}
%\label{eq:mellinED}
H(s,x) = \frac{1}{\Gamma(s)} 
\int_{0}^{\infty} \frac{\Ea(-xt)}{\I_{\alpha}(-t)} t^{s-1}\,dt
= \frac{2}{\Gamma(s)}
\int_{0}^{\infty} \frac{\Ea(-xt)}{\Ea(t)+\Ea(-t)} t^{s-1}\,dt
\end{equation*}
converges in the right plane $\Re(s) > 0$ to a holomorphic function, which may be analytically continued to an entire function of $s$ satisfying
\begin{equation*}
  H(-n,x) = \frac{n!}{\gamma_{n,\alpha}} \E_{n,\alpha}(x),\quad n=0,1,2,\dots
\end{equation*}
\end{theorem}

\begin{proof}
The statement that $H(s,x)$ is convergent for $-1<x<1$ holds by the same reasoning we made in the Bernoulli-Dunkl case. Hence, by Theorem~\ref{teo:MellinAD}, we have $H(-n,x) = n!\,\E_{n,\alpha}(x)/\gamma_{n,\alpha}$, for $n=0,1,2,\dots$
By the same argument as in the Bernoulli-Dunkl case, we can write $A(t)$ as a geometric series as 
\[
A(t) = \frac{1}{\I_{\alpha}(t)} 
= \frac{2}{\Ea(t)}
\frac{1}{1 + \frac{\Ea(-t)}{\Ea(t)}} 
=  \frac{2}{\Ea(t)} \sum_{n=0}^{\infty} \left(\frac{-\Ea(-t)}{\Ea(t)}\right)^n.
\]
The special function $H(s,x)$ is giving (when $-1<x<1$) by
\begin{align*}
H(s,x) & = \frac{1}{\Gamma(s)} \int_{0}^{\infty} \frac{\Ea(-xt)}{\I_{\alpha}(t)} t^{s-1} \,dt 
= \frac{2}{\Gamma(s)} \sum_{n=0}^{\infty}
\int_{0}^{\infty} \frac{\Ea(-xt)}{\Ea(t)} \left(\frac{-\Ea(-t)}{\Ea(t)}\right)^n t^{s-1} \,dt \\
& = \frac{2}{\Gamma(s)} \sum_{n=0}^{\infty}(-1)^n \sigma_1^n 
\int_{0}^{\infty} \frac{\Ea(-xt)}{\Ea(t)} t^{s-1} \,dt,
\end{align*}
which concludes the proof.
\end{proof}

\begin{definition}
\label{def:HD-zeta-euler}
For $x \in (-1, 1)$ and $\Re(s) > 0$, we define the Hurwitz-Dunkl zeta function of Euler type as
\begin{equation}
\label{eq:z-euler-dunkl}
\zeta_{E,\alpha}(s,x) = \frac{1}{\Gamma(s)} \sum_{n=0}^{\infty}
(-1)^n \sigma_1^n  \int_{0}^{\infty} \frac{\Ea(-xt)}{\Ea(t)} t^{s-1} \,dt.
\end{equation}
\end{definition}

Finally, notice that the function $H(s,x)$ may be extended to the entire complex $s$-plane and we have, for $x \in (-1, 1)$ and $\Re(s)>0$,
\begin{equation*}
%\label{eq:HtozetaED}
H(s,x) = 2\zeta_{E,\alpha}(s,x).
\end{equation*}
Hence, we can consider, equivalently, that the analytic extension of
\begin{equation}
\label{eq:HD-zetaHE0}
  \zeta_{E,\alpha}(s,x) 
  =  \frac{1}{\Gamma(s)}
\int_{0}^{\infty} \frac{\Ea(-xt)}{\Ea(t)+\Ea(-t)} t^{s-1}\,dt,
\quad \Re(s) >0
\end{equation}
(which corresponds to~\eqref{eq:z-euler-dunkl}) is
\begin{equation}
\label{eq:HDE-zeta}
  \zeta_{E,\alpha}(s,x) = \frac{1}{2} H(s,x), 
  \quad s \in \C.
\end{equation}

Again, when $\alpha -1/2$ and $x \mapsto 2x-1$ (now by \eqref{eq:EtoE} we make the changes to recover Euler polynomials from Euler-Dunkl polynomials, as we did with the Hurwitz-Dunkl zeta function) we get the function
\begin{equation*}
%\label{eq:z-euler}
\zeta_{E,-1/2}(s,2x-1) = \frac{1}{2^s} \sum_{n=0}^{\infty}
\frac{(-1)^n}{(x+n)^s} = \frac{1}{2^s} \zeta_{E}(s,x),
\end{equation*}
where $\zeta_{E}(s,x)$ is called the Hurwitz-type Euler zeta function (see, for instance,~\cite{SuHu}). Again, $\zeta_{\alpha,E}$ generalized many properties of $\zeta_{E,\alpha}$ by the changes $\alpha = -1/2$ and $x\mapsto 2x-1$ (and sometimes also $t \mapsto t/2$). There is also the recurrence identity 
\[
\zeta_{E,\alpha}(s,x) + \sigma_1(\zeta_{E,\alpha}(s,\cdot))(x) = d_{\alpha}(s,x)
\] 
which generalizes \cite[(2.1)]{SuHu} and \eqref{eq:HDE-zeta} is an integral representation that generalizes \cite[(3.1)]{WiZ}. Moreover, it is easy to prove that the relation of $\zeta_{E,\alpha}$ with Euler-Dunkl polynomials
\begin{equation*}
    \zeta_{E,\alpha}(-n,x) = \frac{1}{2}\E_{n,\alpha}(x)  \frac{n!}{\gamma_{n,\alpha}}
\end{equation*}
holds for all $x \in (-1,1)$ and $n=0,1,2,\dots$, which give, when we recover the classical Hurwitz-type Euler zeta function, the identity (see \cite[(2.7)]{SuHu})
\begin{equation*}
%\label{eq:trad-euler}
\zeta_{E}(-n,x) = \frac{1}{2} E_n(x).
\end{equation*}

%%%%%%%%%%%%%%%%%%%%%%%%%%%%
% The generalized Bernoulli-Dunkl case
%%%%%%%%%%%%%%%%%%%%%%%%%%%%

%-----------
\subsection{The generalized Bernoulli-Dunkl case}
%-----------

In the Bernoulli-Dunkl case we had
\[
 A(t) = \frac{1}{\I_{\alpha+1}(t)} 
 = \frac{t}{\alpha+1} \frac{1}{\Ea(t)}\frac{1}{1 - \frac{\Ea(-t)}{\Ea(t)}}.
\]
For $r$ a positive integer, the generalized Bernoulli-Dunkl polynomials are defined as $A(t)^r E_\alpha(xt) = \sum_{n=0}^{\infty} \B_{n,\alpha}^{(r)}(x)t^n/\cc_{n,\alpha}$, and we have
\[
A(t)^r = \left( \frac{t}{\alpha+1}\frac{1}{\Ea(t)}
\sum_{n=0}^{\infty} \left(\frac{\Ea(-t)}{\Ea(t)}\right)^n\right)^r.
\]

\begin{theorem}
\label{teo:genBD}
Let $\Ea(xt)/(\I_{\alpha+1}(t))^r$ be the generating function of Bernoulli-Dunkl polynomials of order $r=1,2,\dots$ Then for each $x \in (-r,r)$ the integral
\begin{equation}
\label{eq:mellinBDr}
H(s,x) = \frac{1}{\Gamma(s)} \int_{0}^{\infty} 
\frac{\Ea(-xt)}{\left(\I_{\alpha+1}(t)\right)^r} t^{s-1} \,dt
\end{equation} 
converges in the right plane $\Re(s) > r$ to a holomorphic function, which may be analytically continued to an entire function of $s$ satisfying
\begin{equation*}
  H(-n,x) = \frac{n!}{\gamma_{n,\alpha}} \B_{n,\alpha}^{(r)}(x),\quad n=0,1,2,\dots
\end{equation*}
\end{theorem}

The theorem can be easily proved by the same arguments as in Theorem~\ref{teo:MellinAD} and in Subsection~\ref{subsec:BD}. The only thing left to prove is the convergence of $H(s,x)$ in $x \in (-r,r)$.

\begin{proof}[Proof of Theorem~\ref{teo:genBD}]
Let us first analyze the convergence of the integral \eqref{eq:mellinBDr}.
We use the asymptotic behavior of the Kummer confluent hypergeometric function given in \eqref{eq:asymptotic-1f1}, and proceed as in the proof of Lemma~\ref{lem:Mellin-BD}. If $x > 0$, the ``exponential part'' of the integrand of $H(s,x)$ has size $e^{-t(x+r)}$, so the integral converges if $x<r$. Repeating the argument for $x<0$, we get the convergence in $x\in(-r,r)$.

Let us use that, for $r=1,2,\dots$ and $|z|<1$, 
\[
  \left(\sum_{n=0}^{\infty} z^n\right)^r 
  = \frac{1}{(1-z)^r} = \sum_{n=0}^{\infty} \binom{r+n-1}{n} z^n.
\]
Then,
\begin{align*}
A(t)^r &= \left( \frac{t}{\alpha+1}\frac{1}{\Ea(t)}
\sum_{n=0}^{\infty} \left(\frac{\Ea(-t)}{\Ea(t)}\right)^n\right)^r \\
%= \frac{t^r}{(\alpha+1)^r}\frac{1}{\Ea(t)^r} 
%\left(\sum_{n=0}^{\infty} \left(\frac{\Ea(-t)}{\Ea(t)}\right)^n\right)^r \\
&= \frac{t^r}{(\alpha+1)^r}\frac{1}{\Ea(t)^r} 
\sum_{n=0}^{\infty} \binom{r+n-1}{n} \left(\frac{\Ea(-t)}{\Ea(t)}\right)^n.
\end{align*}
Hence,
\begin{align*}
H(s,x) & = \frac{1}{\Gamma(s)} \int_{0}^{\infty} 
\frac{\Ea(-xt)}{\left(\I_{\alpha+1}(t)\right)^r} t^{s-1} \,dt \\
%% & = \frac{1}{\Gamma(s)} \int_{0}^{\infty} \Ea(-xt) \frac{1}{(\alpha+1)^r} 
%%% \frac{1}{\left(\Ea(t)\right)^r} \sum_{n=0}^{\infty} \binom{r+n-1}{n} 
%%% \left(\frac{\Ea(-t)}{\Ea(t)}\right)^n t^{s+r-1} \,dt \\
& = \frac{1}{\Gamma(s)} \frac{1}{(\alpha+1)^r} 
\sum_{n=0}^{\infty} \binom{r+n-1}{n} \int_{0}^{\infty}
\frac{ \Ea(-xt)}{\left(\Ea(t)\right)^r} 
\left(\frac{\Ea(-t)}{\Ea(t)}\right)^n t^{s+r-1} \,dt \\
%% & = \frac{1}{\Gamma(s)} \frac{1}{(\alpha+1)^r}\sum_{n=0}^{\infty} \binom{r+n-1}{n}
%% \int_{0}^{\infty} \tau_{-1}^n \tau_1^{-(n+r)}
%% \left( \Ea(t\cdot) \right)(-x) t^{s} \,dt \\
&  = \frac{\Gamma(s+r)}{\Gamma(s)} \frac{1}{(\alpha+1)^r}
\sum_{n=0}^{\infty} \sigma_1^n 
\Biggl( \binom{r+n-1}{n} \int_{0}^{\infty} 
\frac{\Ea(-xt)}{\left(\Ea(t)\right)^r} t^{s+r-1} \,dt \Biggr),
\end{align*}
which proves the theorem.
\end{proof}

\begin{definition}
\label{def:zetaBDr}
For $\Re(s) > 1$ we define the Hurwitz-Dunkl zeta function of order $r=1,2,3,\dots$ as
\[
\zeta_{\alpha}^{(r)}(s,x) 
 = \frac{1}{\Gamma(s)} \sum_{n=0}^{\infty} \binom{r+n-1}{n} \sigma_1^n 
 \left(\int_{0}^{\infty} \frac{\Ea(-xt)}{\left(\Ea(t)\right)^r} t^{s-1} \,dt \right).
\]

As in \eqref{eq:sumandoHD} (that is, the case $r=1$), we can define the basic term
\begin{equation*}
%\label{eq:sumandoHDr}
d_{\alpha}^{(r)}(s,x) = \frac{1}{\Gamma(s)} \int_{0}^{\infty} 
\frac{\Ea(-xt)}{(\Ea(t))^r} t^{s-1} \,dt
\end{equation*}
and then write
\[
\zeta_{\alpha}^{(r)}(s,x) 
= \sum_{n=0}^{\infty} \binom{r+n-1}{n} \sigma_1^n d_{\alpha}^{(r)}(s,x).
\]
\end{definition}

Notice that, as the function $H(s,x)$ is extended to the entire complex $s$-plane, and for $\Re(s)>0$ we have
\begin{equation*}
%\label{eq:Hzetaar}
H(s,x) = \frac{(s)_r}{(\alpha+1)^r} \,\zeta_{\alpha}^{(r)}(s+r,x)
\quad\text{for }\, \Re(s)>1-r.
\end{equation*}
Hence, we can define the extension of $\zeta_{\alpha}^{(r)}(s,x)$ to the complex plane by using $\zeta_{\alpha}^{(r)}(s+r,x) = (\alpha+1)^r H(s,x)/(s)_r$,
i.e., by taking
\[
  \zeta_{\alpha}^{(r)}(s,x) = (\alpha+1)^r H(s-r,x)/(s-r)_r,
  \qquad -r<x<r,\quad s \in \C,
\]
which generalizes~\eqref{eq:HD-zeta}.
% Notice that when $r=1$, we have the Bernoulli-Dunkl polynomials.

In the case $\alpha=-1/2$, this kind of zeta functions for the classical generalized Bernoulli polynomials has been studied in~\cite[\S\,4.4]{BayChi}; see also \cite[\S\,4.1]{BayChi} for the classical generalized Euler polynomials. 

%%%%%%%%%%%%%%%%%%%%%%%%%%%%
% The generalized Euler-Dunkl case
%%%%%%%%%%%%%%%%%%%%%%%%%%%%

%-----------
\subsection{The generalized Euler-Dunkl case}
%-----------

Again, as we did with the generalized Bernoulli-Dunkl case, by using the generation function of the generalized Euler-Dunkl polynomials we have $A(t) = 1/\I_{\alpha}(t)$ and
\[
A(t)^r = \left( \frac{2}{\Ea(t)}
\sum_{n=0}^{\infty} \left(-\frac{\Ea(-t)}{\Ea(t)}\right)^n\right)^r.
\]

\begin{theorem}
\label{teo:genED}
Let $\Ea(xt)/(\I_{\alpha}(t))^r$ be the generating function of Euler-Dunkl polynomials of order $r=1,2,\dots$ Then for each $x \in (-r,r)$ the integral
\begin{equation*}
%\label{eq:mellinEDr}
H(s,x) = \frac{1}{\Gamma(s)} \int_{0}^{\infty} 
\frac{\Ea(-xt)}{\left(\I_{\alpha}(t)\right)^r} t^{s-1} \,dt
\end{equation*}
converges in the right plane $\Re(s) > 0$ to a holomorphic function which may be analytically continued to an entire function of $s$ satisfying
\begin{equation*}
  H(-n,x) = \frac{n!}{\gamma_{n,\alpha}} \E_{n,\alpha}^{(r)}(x),\quad n=0,1,2,\dots
\end{equation*}
\end{theorem}

\begin{proof}
We only need to notice that
\[
A(t)^{r} = \frac{2^r}{\Ea(t)^r}
\sum_{n=0}^{\infty} (-1)^n \binom{r+n-1}{n} \left(\frac{\Ea(-t)}{\Ea(t)}\right)^n,
\]
and proceed as in Theorem~\ref{teo:genBD}. 
\end{proof}

\begin{definition}
\label{def:zetaEDr}
For $\Re(s) > 0$ we define the Hurwitz-Dunkl zeta function of Euler type and order $r \in \mathbb{N}$ as
\[
\zeta_{E,\alpha}^{(r)}(s,x) = \sum_{n=0}^{\infty}  (-1)^n \binom{r+n-1}{n} \sigma_1^n d_{\alpha}(s,x).
\]
\end{definition}

Finally, as the function $H(s,x)$ is extended to the entire complex $s$-plane, we have
\begin{equation*}
H(s,x) = 2^r \zeta_{E,\alpha}^{(r)}(s,x)
\end{equation*}
for $\Re(s)>0$. Hence, we can define the extension of $\zeta_{E,\alpha}^{(r)}(s,x)$ to the complex plane by 
\[
  \zeta_{E,\alpha}^{(r)}(s,x) = H(s,x)/2^r,
  \qquad -r<x<r,\quad s \in \C.
\]
% Of course, the case $r=1$ led us to the Euler-Dunkl polynomials.

%-----------
\subsection{The generalized Hermite case}
\label{subsec:Her}
%-----------

The classical Hermite polynomials $H_n(x)$ are giving by the generating function $e^{-t^2+2tx}$, and they are orthogonal on the real line with respect to the weight $e^{-x^2}$. 
A well known generalization of these polynomials is the so-called generalized Hermite polynomials of order $\mu > -1/2$, which are orthogonal on the real line with respect to the weight $\omega_{\mu}(x) = |x|^{2\mu}e^{-x^2}$, that is, they are polynomials $\{H_n^{\mu}(x)\}_{n=0}^{\infty}$ satisfying
\begin{equation*}
    \int_{-\infty}^{\infty} H_m^{\mu}(x) H_n^{\mu}(x)\omega_{\mu}(x)\,dx = 0;
\end{equation*}
see, for instance, \cite{AlS}, \cite[Chapters 1 and~5]{Chi} or \cite[p.~380, problem 25]{Szego}.

In \cite{Ros}, Rosenblum shows that these polynomials can be studied in the context of the Dunkl transform on the real line. This is done by means of 
\begin{equation}
\label{eq:Hergen}
e^{-t^2}E_{\mu}(2xt) = \sum_{n=0}^{\infty} H_{n}^{\mu}(x)\frac{t^n}{n!}
\end{equation}
with $\mu = \alpha + 1/2$. Except by a simple change of variable, this is an Appell-Dunkl sequence in the sense of~\eqref{eq:ad}.

For these polynomials, it is easy to find the analytic extension $H(s, x)$ such that, for $n$ a negative integer, the corresponding value is $H_{n}^{\mu}(x)$, except for a multiplicative constant. Due to the factor $e^{-t^2}$, which appears in \eqref{eq:Hergen},  
the extension given in Theorem~\ref{teo:MellinAD} does not present any problem and is valid for $x\in\R$. The same happens with the integral 
\[
  \frac{1}{\Gamma(s)} \int_{0}^{\infty} e^{-t^2}E_{\mu}(-2xt) \,dt,
\]
which is similar to the ones that appear in Theorems~\ref{teo:mellinBD} or~\ref{teo:mellinED}.

That leads us to the following result.

\begin{theorem}
Let $G(-t,x) = e^{-t^2}E_{\mu}(-2xt)$, with $\mu>-1/2$. Then, for $x \in \R$,
\begin{align*}
H(s,x) & = \frac{1}{\Gamma(s)} \int_{0}^{\infty} e^{-t^2}E_{\mu}(-2xt) \,dt \\
& = \frac{\sqrt{\pi}}{2^s\Gamma(\frac{s+1}{2})}
\,{}_1F_1\Big(\frac{s}{2},\mu+\frac{1}{2},x^2\Big) 
- \frac{\sqrt{\pi}}{2^s\Gamma(\frac{s}{2})} \frac{x}{\mu+\frac{1}{2}}
\,{}_1F_1\Big(\frac{s}{2},\mu+\frac{3}{2},x^2\Big)
\end{align*}
is an entire function of $s$ and satisfies $H(-n,x) = H_{n}^{\mu}(x)$ for $n=0,1,2,\dots$
\end{theorem}

In fact, we have $1/\Gamma(\frac{s}{2})=0$ when $s=-2n$, and $1/\Gamma(\frac{s+1}{2})=0$ when $s=-2n-1$ for $n=0,1,2,\dots$ 
Furthermore, $\Gamma(-n+\frac{1}{2}) 
= \frac{(-1)^n \, 2^{2n}(2n)!}{n!}\sqrt{\pi}$, which means
\begin{align*}
H(-2n,x) & = H_{2n}^{\mu}(x) 
= (-1)^n \, \frac{(2n)!}{n!} \,{}_1F_1\Big({-n},\mu+\frac{1}{2},x^2\Big), 
\\
H(-2n-1,x) & = H_{2n+1}^{\mu}(x) 
= (-1)^n \, \frac{(2n+1)!}{n!} 
\frac{x}{\mu+\frac{1}{2}} \,{}_1F_1\Big({-n},\mu+\frac{3}{2},x^2\Big) .
\end{align*}
This, as expected, is the same as \cite[(2.1.1) and (2.2.1)]{Ros}.

%-----------
\section{Properties of the Hurwitz-Dunkl zeta functions} 
\label{sec:prop-hurwitz}
%-----------

The aim of this section is, firstly, to provide generalization of the Riemann zeta function $\zeta(s) = \sum_{n=1}^{\infty} 1/n^s$ and the Euler-type zeta function $\zeta_E(s) = \sum_{n=1}^{\infty} (-1)^{n+1}/n^s$ (also known as Dirichlet eta function $\eta(s)$) in a Dunkl sense through the functions $\zeta_\alpha(s,x)$ and $\zeta_{E,\alpha}(s,x)$, respectively. We also provide a generalization in a Dunkl sense of the analytic continuation of $\zeta(s)$ (and $\zeta_E(s)$), as well as the so-called reflection formula, and other properties concerning our Hurwitz-Dunkl zeta functions $\zeta_\alpha(s,x)$, $\zeta_{E,\alpha}(s,x)$, $\zeta_\alpha(s)$ and $\zeta_{E,\alpha}(s)$. A conection appears here between these functions and the function $Z_\alpha(s) = \sum_{n=1}^{\infty}1/j_{n}^s$ (and also with $Z_{\alpha+1}(s)$), where $j_n$ are the positive zeros of the Bessel function $J_{\alpha}(x)$. We study $Z_\alpha(s)$ in Section~\ref{sec:zetaR}.

In this section we will state the main results. The proofs are rather technical and require several lemmas. We will postpone them to Section~\ref{sec:proof-of-ths}.

%-----------
\subsection{Theorems for $\zeta_\alpha(s,x)$ and  $\zeta_\alpha(s)$}
\label{subsec:teosza}
%-----------

In this section we are going to give another way of expressing the analytic continuation of $\zeta_{\alpha}(s,x)$ for $\Re(s)<1$, and we will give some consequences that involve the zeros of $J_{\alpha+1}(x)$. 

%% LOS CEROS NO SALEN EXPL\'ICITAMENTE EN EL ENUNCIADO DEL TEOREMA,
%% PERO S\'I EN EL LEMA PREVIO PARA PROBARLO.

It is well known (see \cite[Chapter 15]{Wat} or \cite[\S\,10.21]{OlMa-NIST}) that, for any $\alpha>-1$, the zeros of the Bessel function $J_{\alpha}(x)/x^{\alpha}$ can be written as $j_{m,\alpha}$, $m\in \Z \setminus\{0\}$, with $j_{m,\alpha} = -j_{-m,\alpha}$ and $0< j_{m,\alpha} < j_{m+1,\alpha}$, $m\ge 1$.
Moreover, $j_{m,\alpha} \sim (m + \alpha/2 - 1/4)\pi + o(1/m)$ when $m \to \infty$ (see, for instance, \cite[10.21.19] {OlMa-NIST}).

Now, we are interested in the zeros of $J_{\alpha+1}(x)/x^{\alpha+1}$ so, to avoid confusion, we will denote $s_{m,\alpha} = j_{m,\alpha+1}$; in this way, we will often use $s_{m}$ for $s_{m,\alpha}$. Again, with this notation we have $s_{m,\alpha} = -s_{-m,\alpha}$ and $0< s_{m,\alpha} < s_{m+1,\alpha}$, $m\ge 1$, where $i s_{m,\alpha}$, $m \in \Z \setminus \{0\}$, are the zeros of $\I_{\alpha+1}(x)$ (or the zeros of $J_{\alpha+1}(ix)/(ix)^{\alpha+1}$). For $\alpha=-1/2$ we have $s_{m,-1/2}=\pi m$. Let us also note that $\I_{\alpha}(is_{m,\alpha})$ provides a generalization of the sign sequence $(-1)^m$ because $\I_{-1/2}(is_{m,-1/2})=(-1)^m$.

The first result is the following, which is similar to the classical case that can be found, for instance, in~\cite[\S\,12.4, Theorem~12.3]{Ap-ANT}, and the proof follows the same scheme. 
However, we have now the functions $E_\alpha(t)$, which are much more complicated than $e^t$, and then the proof needs some additional details. In particular, we require the use of our Lemma~\ref{lem:acot-g}. Actually, the zeros $s_{m,\alpha}$ do not explicitely appear in the statement of this theorem, but they will be crucial in the proof of the lemma. 
% As explained above, the proofs are delayed up to Section~\ref{sec:proof-of-ths}.

\begin{theorem}
\label{teo:Isz}
Let $x\in (-1,1)$ and define
\begin{equation}
\label{eq:Isz}
I(s,x) 
% = \frac{1}{2\pi i}\int_{C} g(t) t^{s-1}\,dt 
= \frac{1}{2\pi i}\int_C \frac{\Ea(xt)}{\Ea(-t)-\Ea(t)} \,t^{s-1}\,dt,
\end{equation}
where $C$ is the contour shown in Figure~\ref{fig:contour}. 
Then $I(s,x)$ is an entire function of $s$ and satisfies 
\begin{equation}
\label{eq:zetaaI}
\zeta_{\alpha}(s,x) = \Gamma(1-s)I(s,x)\quad \text{ if }\, \Re(s) > 1,
\end{equation}
where $\zeta_{\alpha}(s,x)$ is the function defined in~\eqref{eq:HD-zeta1}.
\end{theorem}

\begin{figure}[ht]
\begin{tikzpicture}
%real negatives line
\draw (-3,0) -- (0,0);
%inner circle and conexion with outer circle
\draw[very thick] (0.7,0) arc (0:165:0.7) -- (-3,0.1736);
\draw (0,0) node[right] {$0$};
\draw (0,0) node {$\bullet$};
\draw[very thick] (0.7,0) arc (360:195:0.7) -- (-3,-0.1736);
%arrows
\draw[thick, ->] (1,0) arc(0:45:1);
\draw[thick, ->] (-1.6,0.4) -- (-2.4,0.4);
\draw[thick, ->] (-2.4,-0.4) -- (-1.6,-0.4);
%nodes
\draw (1,0) node[right] {$C_2$};
\draw (-2.35,-0.8) node[right] {$C_1$};
\draw (-2.35,0.8) node[right] {$C_3$};
\end{tikzpicture}
\caption{The contour $C$ from Theorem~\ref{teo:Isz}.}
\label{fig:contour}
\end{figure}

Taking into account that \eqref{eq:Isz} is valid in the entire $s$-plane, and that $\zeta_{\alpha}(s,x)$ satisfies \eqref{eq:zetaaI} for $\Re(s)>1$, we can define the following analytic continuation for $\zeta_{\alpha}(s,x)$ in the entire $s$-plane:
\begin{equation}
\label{eq:zetaaACI}
    \zeta_{\alpha}(s,x) = \Gamma(1-s) I(s,x),
\end{equation}
valid for $-1<x<1$. Of course, the analytic continuation of a function is unique, so this function $\zeta_{\alpha}(s,x)$ is the same that we defined in~\eqref{eq:HD-zeta}. From this and by Cauchy's residue theorem it is also possible to prove Theorems~\ref{teo:reszeta1} and~\ref{teo:zeta-n}.

\begin{theorem}
\label{teo:reszeta1}
The function $\zeta_{\alpha}(s,x)$ defined in~\eqref{eq:zetaaACI} is analytic for $s \in \C$ except for a simple pole at $s=1$ with residue $\alpha+1$.
\end{theorem}

Taking $\alpha = -1/2$ and $x \mapsto 2x-1$, since $\zeta_{-1/2}(s,2x-1) = \zeta(s,x)/2^s$ we get that $\zeta(s,x)$ has a simple pole at $s=1$ with residue~$1$, which is what happens in the classical case (see \cite[\S\,12.5, Theorem~12.4]{Ap-ANT}).

The next result was already proved in Proposition~\ref{prop:basic-z}, but later we provide another way to show it, this time starting from \eqref{eq:zetaaACI} and using Cauchy's residues theorem:

\begin{theorem}
\label{teo:zeta-n}
The function $\zeta_{\alpha}(s,x)$ defined in~\eqref{eq:zetaaACI} satisfies, for $x \in (-1,1)$,
\begin{equation}
\label{eq:zeta-n}
\zeta_\alpha(-n,x) = -\B_{n+1}(x)\frac{n!\,(\alpha+1)}{\gamma_{n+1,\alpha}},
\qquad n=0,1,2,\dots
\end{equation}
\end{theorem} 

Another classical result in analytic number theory is the so-called Hurwitz formula (see \cite[\S\,12.7, Theorem 12.6]{Ap-ANT} or \cite[25.13.3]{Ap-NIST}), namely
\begin{equation}
\label{eq:H-form}
\zeta(1-s,x) = \frac{\Gamma(s)}{(2\pi)^s} \left(e^{-\pi s i/2} F(x,s)
+ e^{\pi s i/2} F(-x,s)\right),
\end{equation}
%% NOTA: Esta es la notaci\'on habitual, $F(x,s)$, 
%% con los par\'ametros al rev\'es que en $\zeta(s,x)$
where
\begin{equation}
\label{eq:Fxs}
F(x,s) = \sum_{n=1}^{\infty} \frac{e^{2\pi i n x}}{n^s},
\quad \Re(s) > 1,
\end{equation}
is known as the Lerch (or periodic) zeta function (see \cite[\S\,12.7, equation~(9), p.~257]{Ap-ANT} or \cite[25.13.1]{Ap-NIST}).

In the Dunkl context, this formula can be generalized as follows. The convergence of the series~\eqref{eq:Fxs} is clear, but to prove the convergence of the corresponding series $\mathcal{F}(x,s)$, which we will use in the Dunkl context, will require some effort. 
% (of course, it is done in the proof of the theorem, in Section~\ref{sec:proof-of-ths}). 

\begin{theorem}[Hurwitz-Dunkl formula]
\label{teo:HD-form}
Let $\alpha>-1$ and $\{s_m\}_{m=1}^{\infty}$ be the positive zeros of~$J_{\alpha+1}$.
For $\Re(s)>1$, the function
\begin{equation}
\label{eq:HD-formF}
\mathcal{F}(x,s) = \sum_{m=1}^{\infty} 
\frac{\Ea(xi s_m)}{\I_{\alpha}(i s_m)} \frac{1}{s_m^s}
\end{equation}
converges for every $x \in \R$.
Moreover, for $x \in (-1,1)$ and $\Re(s)>1$, 
the Hurwitz-Dunkl zeta function $\zeta_\alpha(s,x)$ satisfies
\begin{equation}
\label{eq:HD-form}
\zeta_{\alpha}(1-s,x) = \frac{\Gamma(s)}{2}
\left(e^{-\pi s i/2} \mathcal{F}(x,s) + e^{\pi s i/2} \mathcal{F}(-x,s)\right).
\end{equation}
\end{theorem}

We call $\mathcal{F}(x,s)$ the Lerch-Dunkl zeta function since it plays a similar role as the Lerch zeta function $F(x,s)$ (but $\mathcal{F}(x,s)$ is not periodic). In fact, when $\alpha = -1/2$ and $x\mapsto 2x-1$, we have $\mathcal{F}(2x-1,s) = \pi^{-s}F(x,s)$, so \eqref{eq:HD-form} becomes~\eqref{eq:H-form}. 

Now, although the identity \eqref{eq:HD-form} is valid only for $x \in (-1,1)$, the right hand side is valid for $x \in \R$, so we can extend the definition of $\zeta_\alpha(1-s,x)$ for $\Re(s)>1$ by taking
\begin{equation}
\label{eq:HD-form-xR}
\zeta_{\alpha}(1-s,x) = \frac{\Gamma(s)}{2}
\left(e^{-\pi s i/2} \mathcal{F}(x,s) + e^{\pi s i/2} \mathcal{F}(-x,s)\right),
\quad x \in \R.
\end{equation}
Replacing $1-s$ by $s$, we can also define, for $\Re(s)<0$,
\[
\zeta_{\alpha}(s,x) = \frac{\Gamma(1-s)}{2}
\left({-}ie^{\pi s i/2} \mathcal{F}(x,1-s) + i e^{-\pi s i/2} \mathcal{F}(-x,1-s)\right),
\quad x \in \R.
\]

The Hurwitz-Dunkl formula gives us an expression for $\zeta_\alpha(s,x)$ and $x \in \R$ free of the intricate integrals. With that, we can easily prove a ``reflection formula'' (but in this case that isn't a suitable name) for $\zeta_\alpha(s)$ in a Dunkl sense that can be seen as a generalization of the reflection formula for~$\zeta(s)$. 

Using the notation $\zeta(s) = \zeta(s,1)$, the reflection formulas of the classical zeta function (also known as ``Riemann's functional equation'')
\begin{align*}
\zeta(1-s) = 2(2\pi)^{-s}\Gamma(s) \cos{\left(\frac{\pi s}{2}\right)}\zeta(s),
\qquad s \in \C,\\
\zeta(s) = 2(2\pi)^{s-1}\Gamma(1-s) \sin{\left(\frac{\pi s}{2}\right)}\zeta(1-s),
\qquad s \in \C,
\end{align*}
can be proved by taking $x=1$ in the Hurwitz formula \eqref{eq:H-form} 
(see, for instance, \cite[\S\,12.8, Theorem~12.7]{Ap-ANT}); 
for the first formula, the result is clear for $\Re(s)>1$, and is then valid for $s \in \C$ by analytic continuation. Actually, many properties of $\zeta(s)$ and $\zeta(s,x)$ can be seen as consequences of~\eqref{eq:H-form}. 

In our case, taking $x= \pm1$ in \eqref{eq:HD-formF}, we get $\mathcal{F}(\pm1,s) = \sum_{m=1}^{\infty} {1}/{s_m^s}$, 
% (that is convergent for $\Re(s) > 1$),
since $\Ea(\pm i s_m) = \I_{\alpha}(i s_m)$.
Thus, we can define, for $\Re(s)>1$,
\begin{equation}
\label{eq:zetaa1-s}
  \zeta_{\alpha}(1-s) = \zeta_{\alpha}(1-s,1),
\end{equation}
where $\zeta_{\alpha}(1-s,1)$ is given in \eqref{eq:HD-form-xR}; of course, the same can be done for $\zeta_{\alpha}(s)$ with $\Re(s)<0$ (with this notation, $\zeta_{-1/2}(s) = \zeta_{-1/2}(s,1) = \zeta(s,1)/2^s = \zeta(s)/2^s$). Then, we have the following:

\begin{theorem}
\label{teo:ref-form}
Let $\alpha>-1$ and $\{s_m\}_{m=1}^{\infty}$ be the positive zeros of~$J_{\alpha+1}$.
For $\Re(s)>1$ we have
\begin{equation}
\label{eq:ZD-cos}
    \zeta_{\alpha}(1-s) = \Gamma(s) \cos{\left(\frac{\pi s}{2}\right)} 
    \sum_{m=1}^{\infty} \frac{1}{s_m^s},
\end{equation}
or equivalently, for $\Re(s)<0$,
\begin{equation}
\label{eq:ZD-sin}
    \zeta_{\alpha}(s) = \Gamma(1-s) \sin\left(\frac{\pi s}{2}\right) 
    \sum_{m=1}^{\infty} \frac{1}{s_m^{1-s}}.
\end{equation}
\end{theorem}

When $\alpha=-1/2$, we get $\sum_{m=1}^{\infty}1/s_m^s = \pi^{-s}\sum_{m=1}^{\infty}1/m^s = \pi^{-s}\zeta(s)$ so, in fact, $\sum_{m=1}^{\infty}1/s_m^s$ is playing the role of $\zeta(s)$. Hence, when $\alpha = -1/2$, Theorem~\ref{teo:ref-form} provides a generalization in a Dunkl sense of classical reflection formulas. However, there is an important difference if we compare Theorem~\ref{teo:ref-form} with the classical case: the sums $\sum_{m=1}^{\infty}$ in \eqref{eq:ZD-cos} and \eqref{eq:ZD-sin} are not the functions $\zeta_{\alpha}(1-s)$ and $\zeta_{\alpha}(s)$, respectively.

Finally, by taking $s=n+1$ in \eqref{eq:HD-form-xR} we have
\begin{equation*}
\zeta_{\alpha}(-n,x) = \frac{n!}{2}
\left(e^{-\pi (n+1) i/2} \mathcal{F}(x,n+1) + e^{\pi (n+1) i/2} \mathcal{F}(-x,n+1)\right),
\quad x \in \R;
\end{equation*}
and, on the other hand, by Theorem~\ref{teo:zeta-n}, 
\begin{equation*}
\zeta_{\alpha}(-n,x) = -\B_{n+1}(x)\frac{n!\,(\alpha+1)}{\gamma_{n+1,\alpha}}, \quad x \in (-1,1).
\end{equation*}
Then, for $x \in (-1,1)$,
\[
-\B_{n+1}(x)\frac{(\alpha+1)}{\gamma_{n+1,\alpha}}
= \frac{1}{2}
\left(e^{-\pi (n+1) i/2} \mathcal{F}(x,n+1) + e^{\pi (n+1) i/2} \mathcal{F}(-x,n+1)\right).
\]

Since the above function is a polynomial, the limit as $x\to 1^{-}$ exists and we have
\[
-\B_{n+1}(1)\frac{(\alpha+1)}{\gamma_{n+1,\alpha}}  = \cos{\left(\frac{\pi (n+1)}{2}\right)} 
   \sum_{m=1}^{\infty} \frac{1}{s_m^{n+1}}.
\]
Letting $n=2k-1$, we have the following:

\begin{corollary}
\label{cor:sum-sm}
Let $\alpha>-1$ and $\{s_m\}_{m=1}^{\infty}$ be the positive zeros of~$J_{\alpha+1}$.
Then,
\begin{equation*}
%\label{eq:sum-bessel}
\sum_{m=1}^{\infty} \frac{1}{s_m^{2k}} 
= \frac{\B_{2k}(1)(-1)^{k+1}}{2^{2k}k!\,(\alpha+2)_{k-1}},
\quad k=1,2,3,\dots
\end{equation*}
\end{corollary}

The previous expression for $\sum_{m=1}^{\infty} {1}/s_m^{2k}$ in terms of $\B_{2k}(1)$ was proved in \cite[Theorem 4.1]{CDPV} by other methods. 

\begin{corollary}
The function $\zeta_{\alpha}(s)$ defined in Theorem~\ref{teo:ref-form} satisfies
\begin{equation}
\label{eq:zeta-Dunkl-n}
\zeta_\alpha(-n) = -\B_{n+1}(1)\frac{n!\,(\alpha+1)}{\gamma_{n+1,\alpha}},
\quad n=1,2,\dots
\end{equation}
\end{corollary}

\begin{proof}
Taking $s=n=2k$, $k=1,2,\dots$, in \eqref{eq:ZD-cos} and using Corollary~\ref{cor:sum-sm} we get \eqref{eq:zeta-Dunkl-n} for $n$ odd. Taking $s=n=2k-1$, $k=1,2,\dots$ in \eqref{eq:ZD-cos} we then get $\cos{(\pi s/2)}=0$ and hence $\zeta_\alpha(1-n)$. Since $-\B_{n+1}(1)=0$ for $n$ even, this completes the proof. 
\end{proof}

%-----------
\subsection{Theorems for $\zeta_{E,\alpha}(s,x)$ and $\zeta_{E,\alpha}(s)$}
\label{subsec:teoszEa}
%-----------

Here, we are going state some results for $\zeta_E(s,x)$, that will be similar to the results for $\zeta_E(s,x)$ in Subsection~\ref{subsec:teosza}.
Let us recall that $j_m = j_{m,\alpha}$, $m\in \Z\setminus\{0\}$, are the zeros of the Bessel function $J_{\alpha}(x)/x^{\alpha}$, and that they can be ordered so that $j_{m,\alpha} = -j_{-m,\alpha}$ and $0< j_{m,\alpha} < j_{m+1,\alpha}$, $m\ge 1$. Moreover, $i j_{m,\alpha}$, $m \in \Z \setminus \{0\}$, are the zeros of $\I_{\alpha}(x)$ and for $\alpha=-1/2$, $j_{m,-1/2}$ are the zeros of $\I_{-1/2}(it)$, namely the zeros of the cosine. Hence, $ j_{m,-1/2} = (m - 1/2) \pi$ for $m \ge 1$.

We begin with a result that is similar to Theorem~\ref{teo:Isz}:

\begin{theorem}
\label{teo:IEsz}
Let $x\in(-1,1)$ and 
\[
I_E(s,x) = \frac{1}{2\pi i}\int_{C} h(t) t^{s-1}\,dt 
= \frac{1}{2\pi i}\int_C \frac{\Ea(xt)}{\Ea(-t)+\Ea(t)} \,t^{s-1}\,dt,
\]
where $C$ is again the contour shown in Figure~\ref{fig:contour} of Theorem~\ref{teo:Isz}.
Then $I_E(s,x)$ is an entire function of $s$ and satisfies 
\begin{equation*}
%\label{eq:zetaEaIe}
\zeta_{E,\alpha}(s,x) = \Gamma(1-s)I_E(s,x) \quad \text{ if } \Re(s)>0,
\end{equation*}
where $\zeta_{E,\alpha}(s,x)$ is the function defined in~\eqref{eq:HD-zetaHE0}.
\end{theorem}

Of course, this theorem again allows us to give the analytic extension for $\zeta_{E,\alpha}(s,x)$ to the entire $s$-plane, valid for $-1<x<1$. 

Now, the ``Hurwitz-Dunkl formula of Euler type'' is the following:

\begin{theorem}
\label{teo:HD-form-E}
Let $\alpha>-1$ and $\{j_m\}_{m=1}^{\infty}$ be the positive zeros of~$J_\alpha$. 
For $\Re(s)>1$, the function
\[
\mathcal{F}_E(x,s) = \sum_{m=1}^{\infty} 
\frac{\Ea( i j_m x)}{\I_{\alpha+1}(i j_m)} \frac{1}{j_m^{s+1}}
\]
converges for every $x \in \R$.
Moreover, for $x \in (-1,1)$ and $\Re(s)>1$ we have
\begin{equation}
\label{eq:form-HD-Euler}
  \zeta_{E,\alpha}(1-s,x) = -(\alpha+1)\Gamma(s)
  \left( e^{-\frac{\pi i}{2}(s+1)} \mathcal{F}_E(x,s) 
  + e^{\frac{\pi i}{2}(s+1)} \mathcal{F}_E(-x,s)\right).
\end{equation}
\end{theorem}

In the particular case $\alpha=-1/2$, we get $j_m = (m - 1/2)\pi$ for $m \ge 1$. Also, $J_{1/2}(x) = \sqrt{2/(\pi x)}\sin{(x)}$, which leads to $\I_{1/2}(ij_m) = (-1)^{m+1}/j_m$. That means, when taking $x\mapsto 2x-1$, 
we have that
\[
\mathcal{F}_E(2x-1,s) 
= \frac{2^s}{\pi^s}i\sum_{m=1}^{\infty} \frac{e^{(2m-1)i\pi x}}{(2m - 1)^{s}} 
= \frac{2^s}{\pi^s} i \ell_E(s,x),
\]
where the notation $\ell_E(s,x)$ for the above series has already been  used in \cite[(7.1)]{SuHu}. Furthermore, with these changes, and noticing that $\mathcal{F}_E(-x,s) = -\mathcal{F}_E(1-x,s)$, \eqref{eq:form-HD-Euler} transforms into
\[
\frac{1}{2^{1-s}}\zeta_E(1-s,x) = \zeta_{E,-1/2}(1-s,2x-1) 
= \frac{2^s\Gamma(s)}{\pi^s} 
\left(e^{-\frac{i\pi s}{2}}\ell_E(s,x) - e^{\frac{i\pi s}{2}}\ell_E(s,1-x)\right),
\]
which is just~\cite[(7.2)]{SuHu}. 


As in the case of $\zeta_\alpha$, we can use \eqref{eq:form-HD-Euler} to define $\zeta_\alpha(1-s,x)$ for $\Re(s)>1$ and $x \in \R$, as well as $\zeta_\alpha(s,x)$ for $\Re(s)<0$ and $x \in \R$.
In particular, taking $x=1$ and defining $\zeta_{E,\alpha}(s) = \zeta_{E,\alpha}(s,1)$, we have the following:

\begin{theorem}
\label{teo:ext-HD-E}
Let $\alpha>-1$ and $\{j_m\}_{m=1}^{\infty}$ be the positive zeros of~$J_{\alpha}$.
For $\Re(s) > 1$ we have
\begin{equation}
\label{eq:ext-HD-E}
\zeta_{E,\alpha}(1-s) = -\Gamma(s)\cos{\left(\frac{\pi s}{2}\right)}
\sum_{m=1}^{\infty} \frac{1}{j_m^s}
\end{equation}
or equivalently, for $\Re(s) < 0$, 
\begin{equation*}
\zeta_{E,\alpha}(s) = -\Gamma(1-s)\sin{\left(\frac{\pi s}{2}\right)}
\sum_{m=1}^{\infty} \frac{1}{j_m^{1-s}}.
\end{equation*}
\end{theorem}

When $\alpha=-1/2$, \eqref{eq:ext-HD-E} transforms into (see \cite[(7.4)]{SuHu})
\[
\zeta_{E}(1-s) 
= -2\pi^{-s}\Gamma(s)\cos{\left(\frac{\pi s}{2}\right)}
\sum_{m=1}^{\infty}\frac{1}{(2m-1)^s}.
\]

Finally, the equivalent result of Corollary~\ref{cor:sum-sm} for $\sum_{m=1}^{\infty} 1/j_m^{2k}$ is the following:

\begin{corollary}
Let $\alpha>-1$ and $\{j_m\}_{m=1}^{\infty}$ be the positive zeros of~$J_\alpha$. 
Then,
\begin{equation*}
%\label{eq:sum-bessel-j}
\sum_{m=1}^{\infty} \frac{1}{j_m^{2k}} 
= \frac{\E_{2k-1}(1)(-1)^{k+1}}{2^{2k}(k-1)!\,(\alpha+1)_{k}},
\qquad k=1,2,3,\dots
\end{equation*}
\end{corollary}

This can be easily proved from~\eqref{eq:ext-HD-E}, as in Subsection~\ref{subsec:teosza}.
Note that this result was also proved in~\cite{DPV} in a different way.

\begin{corollary}
The function $\zeta_{E,\alpha}(s)$ defined in Theorem~\ref{teo:ext-HD-E} satisfies
\begin{equation*}
%\label{eq:zeta-Dunkl-E-n}
\zeta_{E,\alpha}(-n) = \frac{1}{2}\E_{n}(1)\frac{n!}{\gamma_{n,\alpha}},
\quad n=1,2,\dots
\end{equation*}
\end{corollary}

%-----------
\section{Analytic continuation of $\zeta_\alpha(s)$ and  $\zeta_{E,\alpha}(s)$}
\label{sec:zetaR}
%-----------

Finally, let us define, for $\alpha>-1$,
\begin{equation}
\label{eq:zetaR}
  Z_{\alpha}(s) = \sum_{m=1}^{\infty} \frac{1}{j_m^{s}},
  \qquad \Re(s) > 1.
\end{equation}
Of course, in a like manner we get
\begin{equation*}
%\label{eq:zetaR1}
  Z_{\alpha+1}(s) = \sum_{m=1}^{\infty} \frac{1}{s_m^{s}},
  \qquad \Re(s) > 1,
\end{equation*}
hence, $Z_{\alpha}(s)$ is related with $\zeta_{E,\alpha}(s)$ and $Z_{\alpha+1}(s)$ with $\zeta_{\alpha}(s)$. This function is similar to the classical Riemann zeta function $\sum_{m=1}^{\infty} 1/m^s$ where the positive zeros $\{\pi m\}_{m=1}^{\infty}$ of the sine have been changed by the zeros of the positive zeros of a Bessel function. Then, we will call $Z_{\alpha+1}(s)$ the ``Riemann-Bessel zeta function''.

In his thesis \cite{Haw}, Hawkins provides an analytic continuation of $Z_{\alpha}(s)$. To do so, he first gets easily the analytic continuation for $\Re(s)>0$ by integration by parts, and repeating the process he is able to continue the function to $\Re(s)>-1$. However, he does not go forward by this method and, instead, uses other tools. He ends up proving that there exists an analytic continuation of $Z_{\alpha}(s)$ to the entire $s$-plane with simple poles at $s=1,-1,-3,-5,\dots$ but he didn't get an explicit formula. Due to its simplicity, we now show how to continue $Z_{\alpha}(s)$ to the region $\Re(s)>0$.

\begin{theorem}
\label{teo:zetaRanal0}
The function $Z_{\alpha}(s) - \frac{\pi^{-s}}{s-1}$ extends analytically to the region $\Re(s)>0$.
\end{theorem}

\begin{proof}
We start from \eqref{eq:zetaR}, valid for $\Re(s)>1$, and decompose
\begin{align*}
  Z_{\alpha}(s) - \frac{\pi^{-s}}{s-1} 
  &= \sum_{m=1}^{\infty} \frac{1}{j_m^s} - \int_1^{\infty} \frac{\pi^{-s}\,dx}{x^s}
  = \sum_{m=1}^{\infty} \frac{1}{j_m^s} - \sum_{m=1}^{\infty} \int_m^{m+1} \frac{dx}{(\pi x)^s} \\
  &= \sum_{m=1}^{\infty} \left(\frac{1}{j_m^s} - \int_m^{m+1} \frac{dx}{(\pi x)^s} \right)
  = \sum_{m=1}^{\infty} \int_m^{m+1} \left( \frac{1}{j_m^s} - \frac{1}{(\pi x)^s} \right) dx,
\end{align*}
again valid for $\Re(s)>1$. Now, let us denote
\[
  f_m(s) = \int_m^{m+1} \left( \frac{1}{j_m^s} - \frac{1}{(\pi x)^s} \right) dx.
\]
If we prove that $\sum_{m=1}^{\infty} f_m(s)$ is analytic in $\Re(s)>0$, we will have the analytic extension of $Z_{\alpha}(s) - {\pi^{-s}}/(s-1)$.

Every function $f_m(s)$ is analytic in $\Re(s)>0$, so it is enough to see that the series converges uniformly on compacts in that region.
% Si una sucesi\'on de funciones anal\'iticas converge uniformemente sobre compactos 
% a otra funci\'on, esta \'ultima funci\'on tambi\'en es anal\'itica.
Now, let us recall that the zeros $\{j_m\}_{m=1}^\infty$ of $J_{\alpha}(t)$ satisfy 
$j_m \sim (m + \alpha/2 - 1/4)\pi + o(1/m)$ (see \cite[10.21.19] {OlMa-NIST}),
so $\pi m-c \le j_{m} \le \pi m+c$ for a positive constant $c$ independent of~$m$.
Then, because
\[
  \int_{j_m}^{\pi x} u^{-s-1}\,du 
%  = -\frac{1}{s} u^{-s} \Big|_{u=j_m}^{\pi x} 
  = \frac{1}{s} \left( \frac{1}{j_m^s} - \frac{1}{(\pi x)^s} \right),
\]
we have
\begin{align*}
  &\left| \int_m^{m+1} \left( \frac{1}{j_m^s} - \frac{1}{(\pi x)^s} \right) dx \right|
  = \left| s \int_m^{m+1} \! \int_{j_m}^{\pi x} \frac{du}{u^{s+1}} \,dx \right| \\
  &\qquad\le |s| \int_m^{m+1} \! \int_{\pi m-c_1}^{\pi m+c_2} \frac{du}{|u^{s+1}|} \,dx 
  = |s| \int_m^{m+1} \! \int_{\pi m-c_1}^{\pi m+c_2} \frac{du}{u^{1+\Re(s)}} \,dx \\
%  \le |s| \int_m^{m+1} \! \int_m^{m+1} \frac{du}{u^{1+\Re(s)}} \,dx \\
  &\qquad\le \frac{|s|}{(s_m-c_2-c_1)^{1+\Re(s)}} \int_m^{m+1} \int_{\pi m-c_1}^{\pi m+c_2} du \,dx
  = \frac{(c_2+c_1) |s|}{(j_m-c_2-c_1)^{1+\Re(s)}}.
\end{align*}
Consequently,
\[
  \sum_{m=1}^{\infty} \left| \int_m^{m+1} \left( \frac{1}{j_m^s} - \frac{1}{(\pi x)^s} \right) dx \right|
  \le \sum_{m=1}^{\infty} \frac{C |s|}{j_m^{1+\Re(s)}}
  < \infty \quad\text{for }\Re(s) > 0,
\]
and the Weierstrass M-test ensures the uniform convergence on compacts in $\Re(s)>0$.
\end{proof}

Using the analytic continuation in the entire $s$-plane \cite{Haw}, some of the above identities can also be analytically continued to the entire $s$-plane; see also~\cite{Sto}. Since Hawkins didn't get any explicit formula for $Z_{\alpha}(s)$, we won't get an explicit formula for $\zeta_\alpha(s)$ either. Furthermore, in this way we do not obtain $\zeta_{\alpha}(s)$ as our $\zeta_{\alpha}(s,1)$, because $\zeta_{\alpha}(s,x)$ does not exist for $x=1$ in the half-plane $\Re(s) \ge 1$ (see Definition~\ref{def:HD-zeta}). In the same way, contrary to what happen in the classical case, we do not have $\zeta_{\alpha}(s,1) = Z_{\alpha+1}(s)$ for $s\ge 1$. 

Hawkins also computed the residues of $Z_{\alpha}(s)$ at $s=1,-1,-3,\dots$ and the values of $Z_{\alpha}(-2k)$ for $k=0,1,2,\dots$ They are given by (see \cite[Theorem~3.5]{Haw})
\begin{align*}
\Res_{s=-2k-1} ( Z_{\alpha}(s) ) = \frac{(-1)^{k+1}}{\pi}c_{2k},
& &
 Z_{\alpha}(-2k) = \frac{(-1)^{k}}{2}c_{2k-1},
\end{align*}
where $c_k := c_{k,\alpha}$ are given by the identity
\begin{equation}
\label{eq:coef-haw}
\left(\sum_{k=0}^{\infty} \frac{(-1)^k}{(2x)^k} (\alpha,k)\right) 
\left(\sum_{k=0}^{\infty} \frac{c_k}{x^{k+2}} \right)
=
\sum_{k=0}^{\infty} \frac{(-1)^k}{(2x)^{k+1}}(\alpha,k), 
\end{equation}
with
$(\alpha,k) = \Gamma(\alpha+k+1/2) / (k!\,\Gamma(\alpha-k+1/2))$
%\begin{equation*}
%(\alpha,k) = \frac{1}{k!} \frac{\Gamma(\alpha+k+1/2)}{\Gamma(\alpha-k+1/2)}
%\end{equation*}
(see \cite[Lemma~3.4]{Haw}). To extend $\zeta_\alpha(s)$, we need to change $\alpha\mapsto \alpha+1$ in order to correspond the coefficients $c_k$ with our $Z_{\alpha+1}(s)$. Finally, Hawkins proved that $c_k$ are polynomials of $\alpha$ which vanish at $\alpha=-1/2$ and $\alpha=1/2$ (see, for instance, \cite[Proposition~4.2]{Haw}).

With all this, we can now study the poles of $\zeta_\alpha(s)$ with $s\in \C$ and $-1/2 \not = \alpha>-1$.

\begin{theorem}
\label{teo:ZD-ext}
We get
\begin{equation}
\label{eq:ZD-cos-ext}
\zeta_{\alpha}(1-s) = \Gamma(s) \cos{\left(\frac{\pi s}{2}\right)} Z_{\alpha+1}(s), \qquad s \in \C,
\end{equation}
or equivalently
\begin{equation}
\label{eq:ZD-sin-ext}
\zeta_{\alpha}(s) = \Gamma(1-s) \sin{\left(\frac{\pi s}{2}\right)}Z_{\alpha+1}(1-s), \qquad s \in \C.
\end{equation}
In particular, $\zeta_\alpha(s)$ can be analytically continued to the entire $s$-plane with simple poles at $s = n = 1,2,3,\dots$ (for $\alpha\not = -1/2$) whose residues are equal to $d_{n-1}/(2n!)$, where $d_n := d_{n,\alpha} = c_{n,\alpha+1}$ with the notation of~\eqref{eq:coef-haw}. Moreover, $\zeta_\alpha(0) = -1/2$.
\end{theorem}

\begin{proof}
We can analytically continue equations \eqref{eq:ZD-cos} and \eqref{eq:ZD-sin} by considering the analytic continuation of $Z_{\alpha+1}(s)$. From the equation \eqref{eq:ZD-cos-ext}
the only possible poles are the ones of $\Gamma(s)$, at $s=0,-1,-2,\dots$, and the ones of $Z_{\alpha+1}(s)$, at $s=1,-1,-3,\dots$ It is straighforward to see that when $s=-2k=0,-2,-4,\dots$ we get $\cos(-\pi k) Z_{\alpha+1}(-2k)\not= 0$ and $\Gamma(-2k)$ has a pole. Hence, at those values there are simple poles. The residue at $s=-2k$ is
\[
\lim_{s\to -2k} (s+2k)\Gamma(s) \cos{\left(\frac{\pi s}{2}\right)} Z_{\alpha+1}(s) = \frac{(-1)^k}{(2k)!} Z_{\alpha+1}(-2k) = \frac{d_{2k-1}}{2(2k)!}.
\] 
When $s=-(2k+1) = -1,-3,-5,\dots$ we prove that $\zeta_\alpha(1-s)$ has a pole at those values by a little trick and the L'H\^opital rule as follows:
\begin{align*}
& \lim_{s\to -2k-1} \Gamma(s) \cos{\left(\frac{\pi s}{2}\right)} 
Z_{\alpha+1}(s) \\
& =\lim_{s\to -2k-1} (s+2k+1)\Gamma(s) (s+2k+1)Z_{\alpha+1}(s) \frac{\cos{\left(\frac{\pi s}{2}\right)}}{(s+2k+1)^2} \\
 & =\Res_{s=-2k-1}(\Gamma(s)) \Res_{s=-2k-1}(Z_{\alpha+1}(s))
 \lim_{s\to -2k-1} \frac{-\pi}{4}\frac{\sin{\left(\frac{\pi s}{2}\right)}}{s+2k+1}.
\end{align*}
Hence, there is a pole at those values. To calculate its residues we compute
\begin{align*}
& \lim_{s\to -2k-1}  (s+2k+1)\Gamma(s) \cos{\left(\frac{\pi s}{2}\right)} Z_{\alpha+1}(s) \\
& =\lim_{s\to -2k-1} (s+2k+1)\Gamma(s) (s+2k+1)Z_{\alpha+1}(s) \frac{\cos{\left(\frac{\pi s}{2}\right)}}{s+2k+1}\\
 & =\Res_{s=-2k-1}(\Gamma(s)) \Res_{s=-2k-1}(Z_{\alpha+1}(s))
 \lim_{s\to -2k-1} \frac{-\pi}{2}\sin{\left(\frac{\pi s}{2}\right)} \\
 & = \frac{(-1)^{k+1}}{(2k+1)!}\frac{\pi}{2}\Res_{s=-2k-1}(Z_{\alpha+1}(s)) = \frac{d_{2k}}{2(2k+1)!}.
\end{align*}
Finally, we consider the case $s=1$. 
We use $\Res_{s=1}(Z_{\alpha+1}(s)) = 1/\pi$. So,
\begin{align*}
\lim_{s\to 1}  \Gamma(s) \cos{\left(\frac{\pi s}{2}\right)} Z_{\alpha+1}(s) 
& = \lim_{s\to 1} (s-1)Z_{\alpha+1}(s)\Gamma(s)
\frac{\cos{\left(\frac{\pi s}{2}\right)}}{s-1}\\
 & = \Res_{s=1}(Z_{\alpha+1}(s))
 \lim_{s\to 1} \frac{\cos{\left(\frac{\pi s}{2}\right)}}{s-1} 
 = -1/2.
 \qedhere
\end{align*}
\end{proof}

Once we have extended $\zeta_{\alpha}(s)$ to the entire $s$-plane (with simple poles at $s=1,2,3,\dots$), 
we use the continuations of equations \eqref{eq:ZD-cos-ext} and \eqref{eq:ZD-sin-ext}, both valid for $s\in\C$, in order to get
\begin{equation}
\label{eq:rel-ZD-ZR}
  Z_{\alpha+1}(s)Z_{\alpha+1}(1-s) 
  = \frac{\zeta_{\alpha}(1-s)\zeta_{\alpha}(s)}{\Gamma(1-s)\Gamma(s)\sin(\pi s/2)\cos(\pi s/2)}
  = \frac{2}{\pi} \zeta_{\alpha}(1-s)\zeta_{\alpha}(s).
\end{equation}

From that, a simple verification leads us to the following functional equation.

\begin{corollary}
The function
\[
\Phi(s) = \sqrt{\frac{2}{\pi}} \frac{\zeta_\alpha(s)}{Z_{\alpha+1}(s)}
\]
satisfies the functional equation $\Phi(s) = 1/\Phi(1-s)$.
\end{corollary}

Next we study the analytic continuation of $\zeta_{E,\alpha}(s)$ which is rather similar to $\zeta_\alpha(s)$.

\begin{theorem}
We get
\begin{equation}
\label{eq:ZDE-cos-ext}
\zeta_{E,\alpha}(1-s) = -\Gamma(s) \cos{\left(\frac{\pi s}{2}\right)} Z_{\alpha}(s), \qquad s \in \C,
\end{equation}
or equivalently
\begin{equation}
\label{eq:ZDE-sin-ext}
\zeta_{E,\alpha}(s) = -\Gamma(1-s) \sin{\left(\frac{\pi s}{2}\right)} Z_{\alpha}(1-s), \qquad s \in \C.
\end{equation}
In particular, $\zeta_{E,\alpha}(s)$ can be analytically continued to the entire $s$-plane with simple poles at $s = n = 1,2,3,\dots$ (for $\alpha\not = \pm 1/2$) whose residues are equal to $-c_{n-1}/(2n!)$. Moreover, $\zeta_{E,\alpha}(0) = -1/2$.
\end{theorem}

Finally, let us mention that, having $Z_{\alpha}(s)$ defined in $0<\Re(s)<1$ (Theorem~\ref{teo:zetaRanal0}, whose proof provides a convergent series to evaluate $Z_{\alpha}(s)$ in this region), one can wonder where the zeros of these functions are. Hawkins did an analysis of the zeros of $Z_{\alpha}(s)$ and provided some results involving zero free regions \cite[Section~2]{Haw} (see also~\cite{Act}). In addition, many of the graphical or numerical methods for finding the zeros of $\zeta(s)$ in the critical strip (see, for instance, \cite{Ar, ArLu} and the references therein) can be adapted to the case of $Z_{\alpha}(s)$. It is then easy to find zeros of $Z_{\alpha}(s)$ that do not satisfy $\Re(s)=1/2$. However, as far as we know, a further analysis of the zeros of $Z_\alpha(s)$ is yet to be done, but doesn't seem to be straighforward at first glance. Is there a deeper theory behind this problem?

%%%%%%%%%%%%%%%%%%
%% Proof of theorems
%%%%%%%%%%%%%%%%%%

%-----------
\section{Proofs of the results of Section~\ref{sec:prop-hurwitz}}
\label{sec:proof-of-ths}
%-----------

In this section we prove the theorems of Subsections~\ref{subsec:teosza} and~\ref{subsec:teoszEa}. 
In Subsection~\ref{subsec:proofsza}, we begin by proving results concerning the Hurwitz-Dunkl zeta function $\zeta_{\alpha}(s,x)$ stated in Subsection~\ref{subsec:teosza}; in Subsection~\ref{subsec:proofszEa}, and with less details, we prove the corresponding results for $\zeta_{E,\alpha}(s,x)$ stated in Subsection~\ref{subsec:teoszEa}.

%-----------
\subsection{The Dunkl zeta function case}
\label{subsec:proofsza}
%-----------

Our goal is to prove Theorem~\ref{teo:Isz}, and then use it to prove Theorems~\ref{teo:reszeta1}, \ref{teo:HD-form} and~\ref{teo:ref-form}. For that, some preliminary results are needed.

\begin{figure}[b]
    \centering
    \begin{tikzpicture}[scale=0.35]
    %rectangle
    \draw[fill=gray] (-2,-5.2) rectangle (2,5.2);
    %circles
    \draw[fill=white] (0,0) circle (0.5);
    \draw[fill=white] (0,2.5) circle (0.5);
    \draw[fill=white] (0,-2.5) circle (0.5);
    \draw[fill=white] (0,4.5) circle (0.5);
    \draw[fill=white] (0,-4.5) circle (0.5);
    %nodes
    \filldraw[black] (0,0) circle (1pt) ;
    \filldraw[black] (0,2.5) circle (1pt) ;
    \filldraw[black] (0,-2.5) circle (1pt) ;
    \filldraw[black] (0,4.5) circle (1pt) ;
    \filldraw[black] (0,-4.5) circle (1pt) ;
    %lines
    \draw[thick,->] (4,0) -- (2.2,0);
    \draw (4,0) node[right] {$0$};
    \draw[thick,->] (4,2.5) -- (2.2,2.5);
    \draw (4,2.5) node[right] {$is_1$};
    \draw[thick,->] (4,-2.5) -- (2.2,-2.5);
    \draw (4,-2.5) node[right] {$-is_1$};
    \draw[thick,->] (4,4.5) -- (2.2,4.5);
    \draw (4,4.5) node[right] {$is_2$};
    \draw[thick,->] (4,-4.5) -- (2.2,-4.5);
    \draw (4,-4.5) node[right] {$-is_2$};
    \end{tikzpicture}
    \caption{A compact subset of the region $S$ from Lemma~\ref{lem:acot-g}.}
    \label{fig:reg-s}
\end{figure}

\begin{lemma}
\label{lem:acot-g}
Let $\{s_m\}_{m=1}^{\infty}$ be the positive zeros of $J_{\alpha+1}(t)$
and let $S=\C\setminus\{0,\pm is_1,\allowbreak \pm is_2, \dots \}$ denote the region that remains when we remove from the $t$-plane the origin and all zeros of $\I_{\alpha+1}(t)$, as in Figure~\ref{fig:reg-s}. Then for $x\in [-1,1]\setminus \{0\}$, the function 
\[
g(t) = \frac{\Ea(xt)}{\Ea(-t)-\Ea(t)}
\]
is bounded on compact subsets of $S$ and compact subsets of $x\in [-1,1]\setminus \{0\}$.
Furthermore, if $\alpha < 1+1/2$, then for $x=0$ the function $g(t)$ is bounded on compact subsets of~$S$.
\end{lemma}

\begin{proof}
We use arguments similar to those of \cite[\S\,2]{DPV}, and reproduce most of them for the sake of completeness (actually, here it is somewhat simpler because \cite{DPV} uses $\alpha\in\C$ and here we have the standard $\alpha>-1$ of the Dunkl context). 
To get started, let us take a large circle $D=\{z\in\C : |z|=A\}$ of radius $A$ with the condition that none of the points $i s_m$, $m\in \Z\setminus\{0\}$, must lie on~$D$. The poles of $g(t)$ inside $D$ are $i s_m$, with $|s_m| < A$, and all of them are simple.
Now, we prove that the value of $A$ can be chosen arbitrarily large and such that there exists some constant $c > 0$ independent of~$A$ (but depending on~$\alpha$) satisfying
\begin{equation}
\label{eq:cotas2}
    |J_{\alpha}(t)| \geq c\, {e^{\Im(t)}}/{|t|^{1/2}}
\end{equation}
for $t\in D$. For that, we proceed based on what is done in \cite[\S\,15.41, p.~498]{Wat}. First, we denote $H^{(1)}_{\alpha}(t)$ and $H^{(2)}_{\alpha}(t)$ as the Bessel functions of the third kind.  We use the equality
\begin{equation}
\label{eq:JH1H2}
   2J_{\alpha}(t) = H^{(1)}_{\alpha}(t) + H^{(2)}_{\alpha}(t),
\end{equation}
and, in addition, the fact that the Bessel functions of the third kind satisfy the estimates
\begin{align}
\label{eq:estH1}
   H_{\alpha}^{(1)}(t) &= \left(\frac{2}{\pi t}\right)^{1/2}
   e^{i(t-\frac{1}{2}\alpha\pi - \frac{1}{4}\pi)}
   \Big( 1 + \eta_{1,\alpha}(t) \Big),
   \\
\label{eq:estH2}
   H_{\alpha}^{(2)}(t) &= \left(\frac{2}{\pi t}\right)^{1/2}
   e^{-i(t-\frac{1}{2}\alpha\pi - \frac{1}{4}\pi)}
   \Big( 1 + \eta_{2,\alpha}(t) \Big),
\end{align}
were $\eta_{1,\alpha}(t)$ and $\eta_{2,\alpha}(t)$ are functions of order $\calO(1/t)$ for large~$|t|$ (see \cite[\S\,15.4, p.~496]{Wat}). Therefore,
\begin{gather*}
   \frac{1}{2} \left(\frac{2}{\pi |t|}\right)^{1/2} e^{-\Im(t)}
   \leq \left|H_{\alpha}^{(1)}(t)\right|
   \leq 2 \left(\frac{2}{\pi |t|}\right)^{1/2} e^{-\Im(t)},
\\
   \frac{1}{2} \left(\frac{2}{\pi |t|}\right)^{1/2} e^{\Im(t)}
   \leq \left|H_{\alpha}^{(2)}(t)\right|
   \leq 2 \left(\frac{2}{\pi |t|}\right)^{1/2} e^{\Im(t)},
\end{gather*}
for $|t|$ large enough. This, together with~\eqref{eq:JH1H2}, gives
\begin{align*}
   2 |J_\alpha(t)|
   &\geq \frac{1}{2} \left(\frac{2}{\pi |t|}\right)^{1/2} e^{|\Im(t)|}
   - 2 \left(\frac{2}{\pi |t|}\right)^{1/2} e^{-|\Im(t)|}
   \\
   &= \frac{1}{2} \left(\frac{2}{\pi |t|}\right)^{1/2} e^{|\Im(t)|}
   \Big(1 - 4 e^{-2|\Im(t)|}\Big)
\end{align*}
for $|t|$ large enough, which proves~\eqref{eq:cotas2} if $|\Im(t)| \geq 1$. On the two arcs of $D$ with $|\Im(t)| \leq 1$, according to~\eqref{eq:JH1H2}, \eqref{eq:estH1} and~\eqref{eq:estH2}, the problem reduces essentially to get a lower bound for $|\cos(t-\frac{1}{2}\alpha\pi - \frac{1}{4}\pi)|$, which can be done by simply choosing $A$ so that to avoid the zeros of the cosine function. This proves~\eqref{eq:cotas2}.
Furthermore,~\eqref{eq:JH1H2}, \eqref{eq:estH1} and~\eqref{eq:estH2} also give 
\begin{equation*}
   |J_{\alpha} (t)| \leq C \, {e^{|\Im(t)|}}/{|t|^{1/2}}
\end{equation*}
for $|t|$ large enough, with a constant $C > 0$ depending only on~$\alpha$. Therefore, for any compact set $K\subset [-1,1]\setminus\{0\}$ the radius $A$ can be chosen with the additional property that there exists $C > 0$ such that, for any $t \in D$ and any $x \in K$,
\begin{align}
\label{eq:cotasalfa}
   |J_{\alpha} (tx)| &\leq C\, {e^{|\Im(tx)|}}/{|tx|^{1/2}},
   \\
\label{eq:cotasalfa+1}
   |J_{\alpha+1} (tx)| &\leq C\, {e^{|\Im(tx)|}}/{|tx|^{1/2}}.
\end{align}
Using~\eqref{eq:cotas2}, \eqref{eq:cotasalfa} and~\eqref{eq:cotasalfa+1}, we get, for $x\in K$ and $t\in D$,
\begin{align*}
&\left| \frac{\Ea(xt)}{\Ea(-t)-\Ea(t)}\right| 
= \left| \frac{\Ea(tx)\,(\alpha+1)}{\I_{\alpha+1}(t) t}\right|
%     = \left| \frac{(\alpha+1)\left[\I_{\alpha}(tx)
%     + \frac{tx}{2(\alpha+1)}\I_{\alpha+1}(tx)\right]}{\I_{\alpha+1}(t) t}\right| \\
= \left| \frac{(\alpha+1) \I_{\alpha}(tx)}{\I_{\alpha+1}(t) t}
+ \frac{x \I_{\alpha+1}(tx)}{2 \I_{\alpha+1}(t)}\right| \\
&\qquad= \left| \frac{J_{\alpha}(itx) i + J_{\alpha+1}(itx)}{2x^{\alpha} J_{\alpha+1}(it)} \right|
%    \leq \frac{C \frac{e^{\Im(itx)}}{|itx|^{1/2}}(i + \alpha + 1)} 
%    {2 x^{\alpha}c \frac{e^{\Im(it)}}{|it|^{1/2}}}
\leq \tilde{c} \,\frac{e^{|\Im(xit)|}/|xit|^{1/2}}{|x|^{\alpha}e^{|\Im(it)|}/|it|^{1/2}} 
= \tilde{c} \,\frac{e^{(|x|-1)|\Re(t)|}}{|x|^{\alpha-1/2}}
\end{align*}
for some constant $\tilde{c}$ depending only on $\alpha$ and~$K$.
This proves the result for $x \in [-1,1]\setminus\{ 0 \}$. Finally, let us study the particular case $x=0$. As $\Ea(0) = 1$, it follows that
\[
|g(t)| = \left| \frac{\alpha+1}{\I_{\alpha+1}(t)t}\right| 
\leq \tilde{c} \, \frac{|t|^{\alpha -1 -1/2}}{e^{|\Re(t)|}}.
\]
Since we can choose $t$ such as $|t|\to \infty$ and $\Re(t)$ is constant, to ensure that $g(t)$ is bounded at $x=0$ on compact subsets of $S$ we have to consider that $\alpha < 1+1/2$.
\end{proof}

We now have the tools for the next step:

\begin{proof}[Proof of Theorem~\ref{teo:Isz}]
For simplicity, let us denote
\[
  g(t) = \frac{\Ea(xt)}{\Ea(-t)-\Ea(t)}.
\]
The contour $C$ in Figure~\ref{fig:contour} is composed of three parts, $C_1$, $C_2$ and $C_3$. We take $C_2$ as a positively oriented circle of radius $0< c < s_1$ (where $s_1$ is the first zero of $J_{\alpha+1}(x)/x^{\alpha+1}$) about the origin. This avoids $C_2$ passing through a zero of $g(t)$. On the other hand, $C_1$ and $C_3$ are the lower and upper edges of a ``cut'' in the $t$-plane along the negative real axis, traversed as shown in Figure~\ref{fig:contour}.
Then,
\begin{equation}
\label{eq:int-div}
   2\pi i I(s,x) = \left( \int_{C_1} + \int_{C_2}+ \int_{C_3}\right)g(t)t^{s-1}\,dt.
\end{equation}
We consider an arbitrary compact disk $|s|\leq M$ and prove that the integrals along $C_1$ and $C_3$ converge uniformly on every such disk. Since the integrand is an entire function of $s$, this will prove that $I(s,x)$ is entire. 

We have $t = re^{-\pi i}$ on $C_1$, $t=re^{\pi i}$ on $C_3$ (with $r$ varying from $c>0$ to~$\infty$) and $g(t) = g(-r)$. 
Also, let us denote $\sigma = \Re(s)$.
Along $C_1$ and $C_3$, for $r \geq 1$,
\[
|t^{s-1}| = r^{\sigma-1} |e^{\pm\pi i(\sigma-1+iy)}| 
= r^{\sigma-1}e^{\pm\pi y} \leq r^{M-1}e^{\pi M}.
\]
Hence on either $C_1$ or $C_3$, for $r\geq 1$,
\begin{equation*}
    \left|g(t) t^{s-1}\right| \leq r^{M-1}e^{\pi M}\left|g(-r)\right|.
\end{equation*}
Following the proof of Lemma~\ref{lem:acot-g}, we find that $g(-r)$ is bounded by 
\[
  \tilde{c} \, \frac{e^{(|x|-1)r}}{|x|^{\alpha-1/2}}.
\]
% $\tilde{c} {e^{(|x|-1)r}}/{|x|^{\alpha-1/2}}$.
That means $|g(t) t^{s-1}| \leq A r^{M}e^{(|x|-1)|r|}$ for some constant $A$ depending on $M$ and~$x$. Since the integral $\int_c^{\infty} r^{M}e^{(|x|-1)r} \,dr$ converges when $c>0$ and $-1< x < 1$, this shows the convergence along $C_1$ and $C_3$ and hence, $I(s,x)$ is entire.

Now, we compute $I(s,x)$ by \eqref{eq:int-div}, taking into account that $t=ce^{i\theta}$ (with $-\pi \leq \theta \leq \pi$) on $C_2$. Let us take
\begin{align*}
    2\pi i I(s,x) & = \int_{\infty}^{c}r^{s-1}e^{-\pi i s}g(-r)\,dr \\
    &\qquad + \int_{-\pi}^{\pi} c^{s-1} e^{i\theta(s-1)} g(ce^{i\theta})i c e^{i\theta}\,d\theta 
%    \\ &\qquad 
    + \int_{c}^{\infty}r^{s-1}e^{\pi i s}g(-r)\,dr.
\end{align*}
The sum of the integrals along $C_1$ and $C_3$ is equal to
\begin{align*}
   \int_{c}^{\infty} r^{s-1} g(-r) ( e^{\pi i s} - e^{-\pi i s})\,dr 
   & = 2i \sin(s\pi) \int_{c}^{\infty}r^{s-1}g(-r)\,dr \\ 
   & =: 2i \sin(s\pi) I_1(s,c),
\end{align*}
and the integral along $C_2$ is equal to
\[
ic^s\int_{-\pi}^{\pi}e^{i\theta s}g(c e^{i\theta})\,d\theta =: i c^s I_2(s,c).
\]
Dividing by $2i$, we get
\[
\pi I(s,x) = \sin(s\pi) I_1(s,c) + \frac{c^s}{2} I_2(s,c).
\]
If we take $c\to 0$, we notice that $I_1(s,c) \to \Gamma(s)\zeta_{\alpha}(s,x)$ if $\sigma > 1$, where $\zeta_{\alpha}(s,x)$ is the function defined in~\eqref{eq:HD-zeta1}, so it only remains to prove that $I_2(s,c) \to 0$ as $c \to 0$. 

Notice that $g(t)$ is analytic in $|t|<s_1$ except on the simple pole at $t=0$. Hence $g(t)t$ is analytic everywhere on $|t|<s_1$ and so it is bounded here, say $g(t) \leq A/|t|$ for some constant $A > 0$ and $|t|=c>0$. Therefore we have
\[
|I_2(s,c)| \leq \frac{c^{\sigma}}{2} \int_{-\pi}^{\pi} e^{-y\theta} \,\frac{A}{c}\,d\theta 
\leq Ae^{\pi |y|}c^{\sigma-1}.
\]
If $\sigma > 1$ and $c\to 0$, we find $I_2(s,c) \to 0$. In conclusion, for $\sigma > 1$,
\[
\pi I(s,x) = \sin(s\pi) \Gamma(s) \zeta_{\alpha}(s,x),
\]
and finally, using that $\Gamma(s)\Gamma(1-s) = \pi/\sin(\pi s)$, we get~\eqref{eq:zetaaI}.
\end{proof}

Let us give the proof of Theorem~\ref{teo:reszeta1}, where we show that the unique singularity of $\zeta_{\alpha}(s,x)$, such as defined in~\eqref{eq:zetaaACI}, is a simple pole at $s=1$. 

\begin{proof}[Proof of Theorem~\ref{teo:reszeta1}]
Since $I(s,x)$ is entire, the only possible singularities of $\zeta_{\alpha}(s,x)$ are the poles of $\Gamma(1-s)$, that is, the points $s=1,2,3,\dots$ But $\zeta_{\alpha}(s,x)$ is analytic for $s>1$, so $s=1$ is the only possible pole of $\zeta_{\alpha}(s,x)$.

If $s$ is any integer, say $s=n$, the integrand in the contour integral for $I(s,x)$ takes the same values on $C_1$ as on $C_3$, and hence the integrals along $C_1$ and $C_3$ cancel, leaving, by Cauchy's residue theorem,
\[
I(n,x) = \frac{1}{2\pi i}\int_{C_2} \frac{\Ea(xt)}{\Ea(-t)-\Ea(t)} t^{n-1}\,dt = \Res_{t=0} \Big(\frac{\Ea(xt)}{\Ea(-t)-\Ea(t)} t^{n-1}\Big).
\]
In particular, when $s=1$ we have
\[
I(1,x) = \Res_{t=0} \Big(\frac{\Ea(xt)}{\Ea(-t)-\Ea(t)}\Big) 
= - \lim_{t\to 0} \frac{t \Ea(xt)(\alpha+1)}{\I_{\alpha+1}(t)t} = -(\alpha+1).
\]
To find the residue of $\zeta_{\alpha}(s,x)$ at $s=1$ we compute the limit
\[
\lim_{s\to 1}(s-1)\zeta_{\alpha}(s,x) 
= - \lim_{s\to 1}(1-s)\Gamma(1-s) I(s,x) 
= -I(1,x)\lim_{s\to 1}\Gamma(2-s) = \alpha+1.
\]
This proves that $\zeta_{\alpha}(s,x)$ has a simple pole at $s=1$ with residue $\alpha + 1$.
\end{proof}

Now that Theorem~\ref{teo:Isz} is proved, we can obtain the expression of 
$\zeta_\alpha(-n,x)$, for $n=0,1,2,\dots$ related to the Bernoulli-Dunkl polynomials:

\begin{proof}[Proof of Theorem~\ref{teo:zeta-n}]
Evaluating at $s=-n$ in \eqref{eq:Isz} we get $\zeta_\alpha(-n,x) = n!\,I(-n,x)$. Applying Cauchy's residue theorem, we have 
\begin{align*}
I(-n,x) & = \Res_{t=0} \left( \frac{\Ea(xt)}{\Ea(-t)-\Ea(t)}t^{-n-1}\right) 
 = \Res_{t=0} \left( -(\alpha+1)\frac{\Ea(xt)}{\I_{\alpha+1}(t)}t^{-n-2}\right) \\
 & =  -(\alpha+1)\Res_{t=0} \left(t^{-n-2} \sum_{m=0}^{\infty}
 \frac{\B_{m,\alpha}(x)}{\gamma_{m,\alpha}}t^m\right) \\
 & =  -(\alpha+1)\lim_{t\to 0} \left(t^{-n-1} \sum_{m=0}^{\infty} 
 \frac{\B_{m,\alpha}(x)}{\gamma_{m,\alpha}}t^m\right) 
 =  -(\alpha+1) \frac{\B_{n+1,\alpha}(x)}{\gamma_{n+1,\alpha}}.
\qedhere
\end{align*}
\end{proof}

Now we are ready to prove the convergence of the Lerch-Dunkl zeta function $\mathcal{F}(x,s)$ defined in~\eqref{eq:HD-formF}, and the Hurwitz-Dunkl formula:

\begin{proof}[Proof of Theorem~\ref{teo:HD-form}]
We begin by proving the convergence of \eqref{eq:HD-formF}, with $\Re(s)>1$, for $x \in \R$.
For real values of the variable, we have
\[
J_{\alpha}(t)^2 + J_{\alpha+1}(t)^2 = \frac{2}{\pi t}(1+o(1)), \quad t \to \infty,
\]
so $\lim_n s_n^{1/2} |J_\alpha(s_n)| = \sqrt{2/\pi}$, and consequently
\[
  |\I_{\alpha}(is_m)| \sim C s_m^{-\alpha-1/2}, \quad t \to \infty.
\]
Moreover, $|\Ea(xis_m)| \le C|xs_m|^{-\alpha-1/2}$ by~\eqref{eq:asymptotic-1f1}. Then,
\[
  \left| \frac{\Ea(xi s_m)}{\I_{\alpha}(i s_m)} \frac{1}{s_m^s} \right|
  \le C |x|^{-\alpha-1/2} s_m^{-s},
\]
and this guarantees the absolute convergence of~\eqref{eq:HD-formF}.

To prove~\eqref{eq:HD-form}, let us consider the contour integral
\begin{equation}
\label{eq:INsx}
    I_N(s,x) = \frac{1}{2\pi i}\int_{C(N)}\frac{\Ea(xt)}{\Ea(-t)-\Ea(t)} t^{s-1}\,dt,
\end{equation}
where $C(N)$ is the loop shown in Figure~\ref{fig:contour-cN}.
We now denote $\sigma = \Re(s)$.

\begin{figure}[ht]
\centering
\begin{tikzpicture}[scale=0.5]
%inner circle and conexion with outer circle
\draw (1,0) arc (0:170:1) -- (-3,0.1736);
\draw (0,0) node[right] {$0$};
\draw (0,0) node {$\bullet$};
%\draw (1,0) node {$\bullet$};
\draw (1,0) arc (360:190:1) -- (-3,-0.1736);
% outer circle
\draw (3,0) arc (0:177:3);
\draw (3,0) arc (0:90:3);
\draw (3,0) arc (360:183:3);
% nodes
\draw (0,2.5) node[right] {$is_N$};
\draw (0,2.5) node {$\bullet$};
\draw (0,3.5) node[right] {$is_{N+1}$};
\draw (0,3.5) node {$\bullet$};
\draw (0,-2.5) node[right] {$-is_N$};
\draw (0,-2.5) node {$\bullet$};
\draw (0,-3.5) node[right] {$-is_{N+1}$};
\draw (0,-3.5) node {$\bullet$};
%arrows
\draw[ultra thick, ->] (1,0) node[right] {};
\draw[ultra thick, ->] (-2.0,0.1736) -- (-2.01,0.1736);
\draw[ultra thick, ->] (-2.01,-0.1736) -- (-2,-0.1736);
\draw[ultra thick, ->] (3,0) -- (3,-0.01);
\end{tikzpicture}
\caption{The contour $C(N)$ from \eqref{eq:INsx}.}
\label{fig:contour-cN}
\end{figure}

First we prove that $\lim_{N\to \infty} I_N(s,x) = I(s,x)$ if $\sigma < 0$. For this it suffices to show that the integral along the outer circle tends to $0$ as $N\to \infty$.

On the outer circle we have $t=Re^{i\theta}$, $-\pi \leq \theta \leq \pi$, hence
\[
  |t^{s-1}| = |R^{\sigma-1} e^{\pi i(\sigma +iy)}| 
  = R^{\sigma-1}e^{-y \pi}  \leq R^{\sigma-1}e^{\pi |y|}. 
\]
Since the outer circle lies in the set $S$ of Lemma~\ref{lem:acot-g}, the integrand is bounded by $AR^{\sigma-1}e^{\pi |y|}$, where $A$ is the bound for $g(t)$ implied by Lemma~\ref{lem:acot-g}; hence, the integral is bounded by $2\pi AR^{\sigma}e^{\pi |y|}$. This tends to $0$ as $R \to \infty$ if $\sigma < 0$. Therefore, replacing $s$ by $1-s$, we see that
\[
  \lim_{N\to \infty}I_N(1-s,x) = I(1-s,x),\quad \text{ if } \sigma > 1.
\]
Since
\[
  \frac{\Ea(xt)}{\Ea(-t)-\Ea(t)} t^{-s} 
  = -\frac{\Ea(xt)(\alpha+1)}{\I_{\alpha+1}(t)} t^{-s-1},
\]
the poles of $g(t)t^{-s}$ are just the zeros of $\I_{\alpha+1}(t)$, say $i s_m$, $m\in \Z\setminus \{0\}$ (we don't take into account the pole at $t = 0$ because $C(N)$ doesn't contain it). We compute $I_N(1-s,x)$ explicitely by Cauchy's residue theorem. We have
\begin{equation}
\label{eq:exp-res}
I_N(1-s,x) = -\sum_{\substack{m=-N\\m \ne 0}}^{m=N} R(m) 
= -\sum_{\substack{m=-N\\m \ne 0}}^{m=N} 
\Res_{t=i s_m}\left(\frac{\Ea(xt)}{\Ea(-t)-\Ea(t)} t^{-s}\right).
\end{equation}
%where
%\begin{equation*}
% R(m) = \Res_{t=i s_m}\left(\frac{\Ea(xt)}{\Ea(-t)-\Ea(t)} t^{-s}\right).
%\end{equation*}
Now, if $m>0$,
\begin{align*}
  -R(m) & = \lim_{t\to i s_m} (t-is_m)\frac{\Ea(xt)(\alpha+1)}{\I_{\alpha+1}(t)} t^{-s-1} \\
  & = \Ea(xis_m)(is_m)^{-s-1}(\alpha+1)\lim_{t\to i s_m}\frac{(t-is_m)}{\I_{\alpha+1}(t)} \\
  & = \Ea(xis_m)(is_m)^{-s-1}(\alpha+1)\frac{1}{\I'_{\alpha+1}(i s_m)}.
\end{align*}
Now, we compute $\I'_{\alpha+1}(i s_m)$ as follows. First, let us write $\I_\alpha(z) = 2^\alpha \Gamma(\alpha+1){I_{\alpha}(z)}/{z^{\alpha}}$, where $I_\alpha$ is the modified Bessel function of the first kind and order~$\alpha$, see~\cite{Wat, OlMa-NIST}. We will use the identities (see, for instance, \cite[10.29.2]{OlMa-NIST})
% \url{http://dlmf.nist.gov/10.29}
\begin{equation}
\label{eq:Ider+}
  I_\alpha'(z) = I_{\alpha+1}(z) + \frac{\alpha}{z}\,I_\alpha (z)
\end{equation}
and
\begin{equation}
\label{eq:Ider-}
  I_\alpha'(z) = I_{\alpha-1}(z) - \frac{\alpha}{z}\,I_\alpha (z).
\end{equation}
By \eqref{eq:Ider+} we have
\begin{equation}
\label{eq:der-besselI}
  \I_\alpha'(z) = 2^\alpha \Gamma(\alpha+1) \left(\frac{I_\alpha'(z)}{z^\alpha}
  - \alpha \, \frac{I_\alpha(z)}{z^{\alpha+1}}\right)
  = \frac{z}{2(\alpha+1)} \, \I_{\alpha+1}(z),
\end{equation}
and by \eqref{eq:Ider-} (with $\alpha+1$ instead of~$\alpha$) we deduce that
\begin{align*}
  \I_{\alpha+1}'(z) &= 2^{\alpha+1} \Gamma(\alpha+2) \left(\frac{I_{\alpha+1}'(z)}{z^{\alpha+1}}
  - (\alpha+1)\frac{I_{\alpha+1}(z)}{z^{\alpha+2}}\right) \\
  &= 2^{\alpha+1} \Gamma(\alpha+2) \left(\frac{I_{\alpha}(z)}{z^{\alpha+1}}
  - 2(\alpha+1)\frac{I_{\alpha+1}(z)}{z^{\alpha+2}}\right) \\
  &= \frac{2(\alpha+1)}{z}(\I_{\alpha}(z)-\I_{\alpha+1}(z)).
\end{align*}
Hence, $\I_{\alpha+1}'(i s_m) = \frac{2(\alpha+1)}{i s_m}\I_{\alpha}(i s_m)$. Therefore we get, for $m =1, 2, \dots$,
\begin{equation}
\label{eq:res+}
  -R(m) = \frac{1}{2}\frac{\Ea(xis_m)}{\I_{\alpha}(i s_m)}(is_m)^{-s}.
\end{equation}
Analogously, for $m = -1, -2, \dots$, we can compute $-R(m)$ the same way as before, but taking into account that $s_{-m} = -s_m$ and knowing that $\I_{\alpha}(t)$ is an even function of~$t$. In this case, we get
\begin{equation}
\label{eq:res-}
    -R(m) = \frac{1}{2}\frac{\Ea(-xis_m)}{\I_{\alpha}(i s_m)}(-is_m)^{-s}.
\end{equation}
By \eqref{eq:res+} and \eqref{eq:res-} we are able to compute \eqref{eq:exp-res}. Indeed,
\begin{equation*}
  I_N(1-s,x) = \frac{i^{-s}}{2} \sum_{m=1}^{N} \frac{\Ea(xis_m)}{\I_{\alpha}(i s_m)s_m^s} 
  + \frac{(-i)^{-s}}{2} \sum_{m=1}^{N} \frac{\Ea(-xis_m)}{\I_{\alpha}(i s_m)s_m^s}.
\end{equation*}
Writing $i^{-s} = e^{-\pi s/2}$ and $(-i)^{-s} = e^{\pi s/2}$ and taking $N \to \infty$ we get
\begin{equation*}
  I(1-s,x) = \frac{1}{2}\left( e^{-\pi s i/2} \sum_{m=1}^{\infty} 
  \frac{\Ea(x i s_m)}{\I_{\alpha}(i s_m)s_m^s}
  + e^{\pi s i/2} \sum_{m=1}^{\infty} \frac{\Ea(-x i s_m)}{\I_{\alpha}(i s_m)s_m^s}\right).
\end{equation*}
Since $\zeta_{\alpha}(1-s,x) = \Gamma(s) I(1-s,x)$, if we call $\mathcal{F}(x,s) = \sum_{m=1}^{\infty} \frac{\Ea(x i s_m)}{\I_{\alpha}(i s_m)s_m^s}$ we finally get the Hurwitz-Dunkl formula~\eqref{eq:HD-form}.
\end{proof}

The Hurwitz-Dunkl formula gives us an expression for $\zeta_\alpha(s,x)$ free of the intricate integrals. With it, we can easily prove Theorem~\ref{teo:ref-form}.

\begin{proof}[Proof of Theorem~\ref{teo:ref-form}]
Taking $x=1$ in the Hurwitz-Dunkl formula \eqref{eq:HD-form}, we get $\mathcal{F}(1,s) = \sum_{m=1}^{\infty} {1}/{s_m^s}$, since $\Ea(\pm i s_m) = \I_{\alpha}(i s_m)$. Hence, 
\begin{equation*}
\zeta_{\alpha}(1-s) = \zeta_{\alpha}(1-s,1) 
= \frac{\Gamma(s)}{2} \sum_{m=1}^{\infty} 
\frac{1}{s_m^s} \left( e^{-\pi s i/2} + e^{\pi s i/2}\right) 
= \Gamma(s) \cos{\left(\frac{\pi s}{2}\right)} \sum_{m=1}^{\infty} \frac{1}{s_m^s} .
\end{equation*}
Changing $1-s$ for $s$ we get the equivalent expression in terms of $\sin(x)$.
\end{proof}

%-----------
\subsection{The Euler-type Dunkl zeta function case}
\label{subsec:proofszEa}
%-----------

Now we prove the analogous results for $\zeta_{E,\alpha}(s,x)$ and $\zeta_{E,\alpha}(s)$.

\begin{lemma}
\label{lem:acot-h}
Let $\{j_m\}_{m=1}^{\infty}$ be the positive zeros of $J_{\alpha}(t)$ and
let $S=\C\setminus\{0,\allowbreak \pm ij_1,\allowbreak \pm ij_2,\allowbreak \dots \}$.
% (i.e., the region that remains when we remove from the $t$-plane the origin 
% and all zeros of $\I_{\alpha}(t)$). 
Then for $x\in [-1,1]\setminus \{0\}$, the function 
\[
h(t) = \frac{\Ea(xt)}{\Ea(-t)+\Ea(t)}
\]
is bounded on compact subsets of $S$ and compact subsets of $x\in [-1,1]\setminus \{0\}$.
Furthermore, if $\alpha < 1/2$, then for $x=0$ the function $h(t)$ is bounded on compact subsets of~$S$.
\end{lemma}

\begin{proof}
We start again by taking a large circle $D=\{z\in\C : |z|=A\}$ of radius $A$ with the only condition that none of the points $j_m$, $m\in \Z\setminus\{0\}$, must lie on~$D$. The poles of $h(t)$ inside $D$ are $j_m$, with $|j_m| < A$, and all of them are simple.
 For any compact set $K\subset [-1,1]\setminus\{0\}$ the radius $A$ can be chosen with the additional property that there exist $c, C > 0$ such that, for any $t \in D$ and $x \in K$, equations~\eqref{eq:cotas2}, \eqref{eq:cotasalfa} and~\eqref{eq:cotasalfa+1} are satisfied. Hence, we have
\begin{align*}
\left| \frac{\Ea(xt)}{\Ea(-t)+\Ea(t)}\right| 
&= \left| \frac{\Ea(tx)}{2 \I_{\alpha}(t)}\right|
= \left| \frac{\I_{\alpha}(tx)}{2\I_{\alpha}(t)}
+ \frac{xt \I_{\alpha+1}(tx)}{4(\alpha+1) \I_{\alpha}(t)}\right| \\
&= \left| \frac{J_{\alpha}(itx) i + J_{\alpha+1}(itx)}{2x^{\alpha} J_{\alpha}(it)} \right|
%    &\leq \dfrac{C \frac{e^{\Im(itx)}}{|itx|^{1/2}}(i + \alpha + 1)}
%    {2 x^{\alpha}c \frac{e^{\Im(it)}}{|it|^{1/2}}}
\leq \tilde{c}\, \frac{e^{(|x|-1)|\Re(t)|}}{|x|^{\alpha-1/2}}
\end{align*}
for some constant $\tilde{c}$ depending only on $\alpha$ and~$K$.
This proves the result for $x \in [-1,1]\setminus\{ 0 \}$. Finally, we consider the particular case $x=0$. As $\Ea(0) = 1$, it follows that
\[
|h(t)| = \left| \frac{1}{2\I_{\alpha}(t)}\right| 
\leq \tilde{c} \,\frac{|t|^{\alpha + 1/2}}{e^{|\Re(t)|}},
\]
which is bounded, when $|t| \geq A$, if $\alpha < 1/2$. Hence, for $x=0$, $h(t)$ is bounded on $S$ if $\alpha < 1/2$. 
\end{proof}

\begin{proof}[Proof of Theorem~\ref{teo:IEsz}]
The proof is identical to the one of Theorem~\ref{teo:Isz} but using the bound of Lemma~\ref{lem:acot-h} instead.
\end{proof}

\begin{proof}[Proof of Theorem~\ref{teo:HD-form-E}]
The convergence of $\mathcal{F}_E(x,s)$, for $x \in \R$, can be proved as in the case of Theorem~\ref{teo:HD-form}, this time with $|\I_{\alpha+1}(j_m)| \sim C |j_m|^{-\alpha-1-1/2}$, so
\[
  \left| \frac{\Ea( i j_m x)}{\I_{\alpha+1}(i j_m)} \frac{1}{j_m^{s+1}} \right|
  \le C |x|^{-\alpha-1/2} j_m^{-s},
\]
and the converges is again for $\Re(s)>1$.

To prove \eqref{eq:form-HD-Euler}, let us now consider 
\[
    I_N(s,x) = \frac{1}{2\pi i} \int_{C(N)}\frac{\Ea(xt)}{\Ea(-t)+\Ea(t)} t^{s-1}\,dt
\]
with $C(N)$ the loop of Figure~\ref{fig:contour-cN}.
On the outer circle the integrand is bounded by $AR^{\sigma-1}e^{\pi |y|}$, where $A$ is the bound for $h(t)$ implied by Lemma~\ref{lem:acot-h}; hence, the integral is bounded by $2\pi AR^{\sigma}e^{\pi |y|}$. If $\sigma < 0$ the integral $I_N(s,x)\to 0$ along the outer circle of $C(N)$ when $R\to \infty$.
Hence, replacing $s$ for $1-s$, we get $\lim_{N\to \infty} I_N(1-s,x) = I_E(1-s,x)$ for $\sigma > 1$. We compute $I_N(1-s,x)$ by Cauchy's residue theorem. Let $m=1,2,\dots$ We compute the residue at $t=i j_m$ using~\eqref{eq:der-besselI}:
\begin{align*}
    -R(m) & = -\Res_{t=i j_m} \left(\frac{\Ea(xt)}{2\I_{\alpha}(t)}t^{-s}\right) 
    = -\lim_{t \to i j_m}(t - ij_m) \left(\frac{\Ea(xt)}{2\I_{\alpha}(t)}t^{-s}\right) \\
    & = -\frac{\Ea(i j_m x)}{2\I_{\alpha}^{'}(i j_m)}(i j_m)^{-s} 
    = -\frac{(\alpha+1)\Ea(i j_m x)}{\I_{\alpha+1}(i j_m)}(i j_m)^{-s-1}.
\end{align*}
Also, when $m=-1,-2,\dots$, we have
\[
    -R(m) = -\frac{(\alpha+1)\Ea(-i j_m x)}{\I_{\alpha+1}(-i j_m)}(-i j_m)^{-s-1}.
\]
Then,
\begin{equation*}
I_N(1-s,x) = -\sum_{\substack{m=-N\\m \ne 0}}^{N} R(m) 
= -\sum_{\substack{m=-N\\m \ne 0}}^{N} 
\frac{(\alpha+1)\Ea(i j_m x)}{\I_{\alpha+1}(i j_m)}(i j_m)^{-s-1}.
\end{equation*}
Letting $N\to \infty$, we get~\eqref{eq:form-HD-Euler}.
\end{proof}

\subsection*{Acknowledgments.} 
The authors wish to thank the anonymous referees and the editor for their detailed revision of the manuscript, which has allowed us to improve the final version of the paper.

%-----------
\begin{thebibliography}{00}
%-----------

\bibitem{Act}
A. Actor and I. Bender,
The zeta function constructed from the zeros of the Bessel function,
\textit{J. Phys. A} \textbf{29} (1996), 6555--6580.

\bibitem{AlS}
W. A. Al-Salam,
Characterization theorems for orthogonal polynomials,
\textit{Orthogonal Polynomials: Theory and Practice}
(Columbus, OH, 1989, P.~Nevai, ed.), 1--24,
NATO Adv. Sci. Inst. Ser. C Math. Phys. Sci., 294,
Kluwer Acad. Publ., Dordrecht, 1990.

\bibitem{Ap}
T. M. Apostol,
On the Lerch zeta function,
\textit{Pacific J. Math.} \textbf{1} (1951), 161--167;
addendum, \textit{Pacific J. Math.} \textbf{2} (1952), 10.

\bibitem{Ap-ANT}
T. M. Apostol,
\textit{Introduction to Analytic Number Theory},
Springer-Verlag, New York, 1976.

\bibitem{Ap-NIST}
T. M. Apostol,
Zeta and related functions,
\textit{NIST handbook of mathematical functions},
% (edited by F.~W. F.~Olver, D.~W. Lozier, R.~F. Boisvert and C.~W. Clark),
587--599, % Es el Chapter 25
U.S. Dept. Commerce, Washington, DC, 2010.
%% LONG CITE:
% National Institute of Standards and Technology, Washington, DC,
% and Cambridge University Press, Cambridge, 2010.
Available online in \url{http://dlmf.nist.gov/25}

% MR2772105
\bibitem{Ar}
J. Arias de Reyna,
High precision computation of Riemann's zeta function by the Riemann-Siegel formula, I,
\textit{Math. Comp.} \textbf{80} (2011), 995--1009. 

% MR3206394 
\bibitem{ArLu}
J. Arias de Reyna and J. van de Lune, 
On the exact location of the non-trivial zeros of Riemann's zeta function,
\textit{Acta Arith.} \textbf{163} (2014), 215--245.

\bibitem{BayChi}
A. Bayad and J. Chikhi,
Reduction and duality of the generalized Hurwitz-Lerch zetas,
\textit{Fixed Point Theory Appl.} \textbf{2013} (2013), article 82, 14~pp.

\bibitem{BaySim}
A. Bayad and Y. Simsek,
Values of twisted Barnes zeta functions at negative integers,
\textit{Russ. J. Math. Phys.} \textbf{20} (2013), 129--137.

\bibitem{BG}
Y. Ben Cheikh and M. Gaied,
Dunkl-Appell $d$-orthogonal polynomials,
\textit{Integral Transforms Spec. Funct.}
\textbf{18} (2007), 581--597.

\bibitem{BKI}
A. Bouanani, L. Kh\'eriji and M. Ihsen Tounsi,
Characterization of $q$-Dunkl Appell symmetric orthogonal $q$-polynomials,
\textit{Expo. Math.}
\textbf{28} (2010), 325--336.

\bibitem{Chi}
T. S. Chihara,
\textit{An Introduction to Orthogonal Polynomials},
Gordon and Breach, New York, 1978.

\bibitem{CDPV}
\'O. Ciaurri, A. J. Dur\'an, M. P\'erez and J. L. Varona,
{Bernoulli-Dunkl and Apostol-Euler-Dunkl polynomials with applications to series involving zeros of Bessel functions},
\textit{J. Approx. Theory}
\textbf{235} (2018), 20--45.

\bibitem{CMV}
\'O. Ciaurri, J. M\'inguez Ceniceros and J. L. Varona,
{Bernoulli-Dunkl and Euler-Dunkl polynomials and their generalizations},
\textit{Rev. R. Acad. Cienc. Exactas F\'is. Nat. Ser. A Mat. RACSAM}
\textbf{113} (2019), 2853--2876.
% https://doi.org/10.1007/s13398-019-00662-z

%% NO SE CITA
%\bibitem{CPRV}
%\'O. Ciaurri, M. P\'erez, J. M. Reyes and J. L. Varona,
%Mean convergence of Fourier-Dunkl series,
%\textit{J. Math. Anal. Appl.}
%\textbf{372} (2010), 470--485.
%% http://www.unirioja.es/cu/jvarona/downloads/Fourier-Dunkl-series.pdf

%% NO SE CITA
%\bibitem{CV}
%\'O. Ciaurri and J. L. Varona,
%{A Whittaker-Shannon-Kotel'nikov sampling theorem related to the Dunkl transform},
%\textit{Proc. Amer. Math. Soc.}
%\textbf{135} (2007), 2939--2947.

%% NO SE CITA
%\bibitem{Cho}
%F. M. Cholewinski,
%\textit{The finite calculus associated with Bessel functions},
%Contemporary Mathematics, 75, American Mathematical Society, Providence, RI, 1988.

\bibitem{Dil-NIST}
K. Dilcher,
Bernoulli and Euler polynomials,
\textit{NIST handbook of mathematical functions},
% (edited by F.~W. F.~Olver, D.~W. Lozier, R.~F. Boisvert and C.~W. Clark),
587--599, % Es el Chapter 24
U.S. Dept. Commerce, Washington, DC, 2010.
%% LONG CITE:
% National Institute of Standards and Technology, Washington, DC,
% and Cambridge University Press, Cambridge, 2010.
Available online in \url{http://dlmf.nist.gov/24}

\bibitem{DH}
I. H. Dimovski and V. Z. Hristov,
Nonlocal operational calculi for Dunkl operators,
\textit{SIGMA Symmetry Integrability Geom. Methods Appl.}
\textbf{5} (2009), Paper~030, 16~pp.

\bibitem{Du}
C. F. Dunkl,
{Differential-difference operators associated to reflection groups},
\textit{Trans. Amer. Math. Soc.}
\textbf{311} (1989), 167--183.

\bibitem{DPV}
A. J. Dur\'an, M. P\'erez and J. L. Varona,
{Fourier-Dunkl system of the second kind and Euler-Dunkl polynomials},
\textit{J. Approx. Theory}
\textbf{245} (2019), 23--39.
% https://doi.org/10.1016/j.jat.2019.04.007

%% NO SE CITA
%\bibitem{DPV-bessel}
%A. J. Dur\'an, M. P\'erez and J. L. Varona,
%{On the properties of zeros of Bessel series in the real line}, 
%\textit{Integral Transforms Spec. Funct}. 
%\textbf{32} (2021), 912--931.

\bibitem{ExLaMiVa}
J. I. Extremiana Aldana, E. Labarga, J. Mínguez Ceniceros and J. L. Varona, 
Discrete Appell-Dunkl sequences and Bernoulli-Dunkl polynomials of the second kind, 
\textit{J. Math. Anal. Appl.} \textbf{507} (2022), 125832, 20~pp. 

\bibitem{DKP} % Tiene Bernoulli
R. Graham, D. Knuth and O. Patashnik, 
\textit{Concrete mathematics: A foundation for computer science}, 
2nd ed., Addison-Wesley, 1994.

%% NO SE CITA
%\bibitem{Erdelyi}
%A. Erd\'elyi, W. Magnus, F. Oberhettinger and F. Tricomi,
%\textit{Higher Transcendental Functions}, 
%Volume I, McGraw-Hill, New York, 1953.

%% NO SE CITA
%\bibitem{EMOT2}
%A. Erd\'elyi, W. Magnus, F. Oberhettinger and F. G. Tricomi,
%\textit{Tables of integral transforms},
%vol.~II, McGraw-Hill, New York, 1954.

%% NO SE CITA
%\bibitem{Guillera}
%J. Guillera and J. Sondow, 
%Double integrals and infinite products for some classical constants via analytic continuations of Lerch's transcendent, 
%\textit{Ramanujan J.} 
%\textbf{16} (2008), 247--270.

% MR2633293
\bibitem{Haw}
J. H. Hawkins, 
\textit{On a zeta-function associated with Bessel's equation}, 
Ph.D. Thesis, University of Illinois at Urbana-Champaign, 1983. 
% https://www.proquest.com/pagepdf/303279356?accountid=17251

\bibitem{SuHu}
S. Hu, D. Kim and M.-S. Kim, 
Special values and integral representations for the Hurwitz-type Euler zeta functions,
\textit{J. Korean Math. Soc.} 
\textbf{55} (2018), no. 1, 185--210.

%% NO SE CITA
%\bibitem{Jeu}
%M. F. E. de Jeu,
%{The Dunkl transform},
%\textit{Invent. Math.}
%\textbf{113} (1993), 147--162.

%% NO SE CITA
%\bibitem{Le}
%N. N. Lebedev,
%\textit{Special functions and their applications},
%Dover, New York, 1972.

\bibitem{LiLiao}
Z. Li and J. Liao % Zhongkai Li and Jianquan Liao
Harmonic analysis associated with the one-dimensional Dunkl transform,
\textit{Constr. Approx.} \textbf{37} (2013), 233--281.

\bibitem{NRV-2018}
L. M. Navas, F. J. Ruiz and J. L. Varona,
Appell polynomials as values of special functions,
\textit{J. Math. Anal. Appl.}
\textbf{459} (2018), 419--436.

\bibitem{NRV-2019}
L. M. Navas, F. J. Ruiz and J. L. Varona,
A note on Appell sequences, Mellin transforms and Fourier series,
\textit{J. Math. Anal. Appl.}
\textbf{476} (2019), 836--850.

% Primeros Bernoulli de orden alpha (entero positivo)
\bibitem{Nor}
% N. E. N\"orlund, % <--- In the paper, this is the name of the author.
N. E. N{\o}rlund, % <--- The correct name in Danish
M\'emoire sur les polyn\^omes de Bernoulli,
\textit{Acta. Math.}
\textbf{43} (1922), 121--196.

% Primeros Euler de orden alpha (entero positivo y negativo)
\bibitem{Nor-book}
% N. E. N\"orlund, % <--- In the book, this is the name of the author.
N. E. N{\o}rlund, % <--- The correct name in Danish
\textit{Vorlesungen \"uber Differenzenrechnung}, 1st ed.,
Springer-Verlag, Berlin-Heidelberg, 1924.

\bibitem{OlMa-NIST}
F. W. J. Olver and L. C. Maximon,
Bessel Functions,
\textit{NIST handbook of mathematical functions},
% (edited by F.~W. F.~Olver, D.~W. Lozier, R.~F. Boisvert and C.~W. Clark),
215--286, % Es el Chapter 10
U.S. Dept. Commerce, Washington, DC, 2010.
%% LONG CITE:
% National Institute of Standards and Technology, Washington, DC,
% and Cambridge University Press, Cambridge, 2010.
Available online in \url{http://dlmf.nist.gov/10}

\bibitem{Ros}
M. Rosenblum,
{Generalized Hermite polynomials and the Bose-like oscillator calculus},
\textit{Oper. Theory Adv. Appl.}
\textbf{73} (1994), 369--396.

\bibitem{Rosler-LN}
M. R\"osler,
Dunkl operators: theory and applications,
\textit{Orthogonal polynomials and special functions
(Leuven, 2002)}, 93--135,
Lecture Notes in Math., 1817, Springer, Berlin, 2003.

\bibitem{Rosler-TAMS}
M. R\"osler,
A positive radial product formula for the Dunkl kernel,
\textit{Trans. Amer. Math. Soc.}
\textbf{355} (2003), 2413--2438.

%\bibitem{Rudin}
%W. Rudin, 
%\textit{Priciples of Mathematical Analysis}, 3rd ed.,
%International Series in Pure and Applied Mathematics,
%McGraw-Hill, 1976.

%% Director de tesis de Hawkins
\bibitem{Sto}
K. B. Stolarsky,
Singularities of Bessel-zeta functions and Hawkins' polynomials,
\textit{Mathematika}
\textbf{32} (1985), 96--103.

\bibitem{Szego}
G. Szeg\H{o}, 
\textit{Orthogonal Polynomials}, 4th~ed.,
Amer. Math. Soc. Colloquium Publ., Vol.~23, 
American Mathematics Society, New York, 1975.
% La primera edici\'on es de 1939

\bibitem{ThX}
S. Thangavelu and Y. Xu,
Riesz transform and Riesz potentials for Dunkl transform,
\textit{J. Comput. Appl. Math.}
\textbf{199} (2007), 181--195.

\bibitem{Wat}
G. N. Watson,
\textit{A Treatise on the Theory of Bessel Functions}, 2nd ed.,
Cambridge Univ. Press,
Cambridge, 1944.

\bibitem{WiZ}
K. S. Williams and N. Y. Zhang, 
Special values of the Lerch zeta function and the evaluation of certain integrals, 
\textit{Proc. Amer. Math. Soc.} \textbf{119} (1993), no. 1, 35--49.


%-----------
\end{thebibliography}
%-----------

%%%%%%%%%%%%%%%%%%%%%%%%%%%%
%---------------------
\end{document}
%---------------------
%%%%%%%%%%%%%%%%%%%%%%%%%%%%
